% concatenated from dualmot.tex
\documentclass[twoside,11pt,a4paper]{article}
\usepackage{jmlr}


\usepackage[utf8]{inputenc}
\usepackage[T1]{fontenc}


\usepackage{macros}


\usepackage{xcolor}
\definecolor{darkblue}{rgb}{0.0,0.0,0.7}
\definecolor{linkequation}{rgb}{0.0, 0.0, 1.0}


\usepackage[
  linktocpage=true,
  hyperfootnotes=false,
  colorlinks=true,
  citecolor=darkblue,
  urlcolor=darkblue,
  linkcolor=magenta,
  backref=page
]{hyperref}


\renewcommand\backreftwosep{, }
\renewcommand*{\backrefalt}[4]{  \ifcase #1    
  \or (Cited on page:~#2.)  \else (Cited on pages:~#2.)  \fi
}


\newcommand*{\SavedEqref}{}
\let\SavedEqref\eqref
\renewcommand*{\eqref}[1]{  \begingroup
    \hypersetup{
      linkcolor=linkequation,
      linkbordercolor=linkequation,
    }    \SavedEqref{#1}  \endgroup
}

\newcommand*{\SavedRef}{}
\let\SavedRef\ref
\renewcommand*{\ref}[1]{  \begingroup
    \hypersetup{
      linkcolor=linkequation,
      linkbordercolor=linkequation,
    }    \SavedRef{#1}  \endgroup
}


\usepackage[square,numbers]{natbib}


\usepackage[colorinlistoftodos]{todonotes}


\usepackage[symbol*]{footmisc}
\DefineFNsymbolsTM{myfnsymbols}{}

\setfnsymbol{myfnsymbols}

\usepackage{multicol}
\usepackage{tabularx}
\usepackage{authblk}

\begin{document}


\title{A Unified Kantorovich Duality for Multimarginal Optimal Transport}

\author[1]{Yehya Cheryala}
\author[1]{Mokhtar Z. Alaya}
\author[1]{Salim Bouzebda}
\affil[1]{{\footnotesize Université de Technologie de Compiègne,  \authorcr
LMAC (Laboratoire de Mathématiques Appliquées de Compiègne), CS 60 319 - 60 203 Compiègne Cedex
}
\texttt{e-mails}: \texttt{yehya.cheryala@utc.fr}; \texttt{alayaelm@utc.fr}; \texttt{salim.bouzebda@utc.fr}}

\maketitle




\begin{abstract}
We study Kantorovich duality for multimarginal optimal transport (MOT) with bounded
continuous cost functions. The main focus is not only the equality between the
primal and dual values, but also the structure of optimal dual potentials. For
compact metric spaces, we prove that the dual problem admits an optimizer in the
class of mutually $c$-conjugate families. The proof combines a
Fenchel--Rockafellar duality argument with equicontinuity estimates, a
normalization of the dual potentials, and the Arzelà--Ascoli theorem. We then consider the non-compact case, where the marginal spaces are Polish. We
first recover the Kantorovich duality identity by a truncation and tightness
argument, using the boundedness of the cost and Prokhorov compactness. For dual
attainment, we rely on the geometry of optimal supports. Under a natural
support-splitting condition, we obtain a bounded Borel measurable dual optimizer
in the canonical class of mutually $c$-conjugate families. Thus, the paper identifies mutually $c$-conjugate potentials as natural
representatives of optimal dual families in the multimarginal setting. These
results provide a structural basis for further work on stability, empirical
multimarginal transport, and statistical limit theory.
\end{abstract}


\begin{keywords}
Multimarginal optimal transport; Kantorovich duality; mutually \(c\)-conjugate potentials; \(c\)-splitting sets;
dual attainment.
\end{keywords}




\section{Introduction}\label{sec:intro}

This introduction places the present work within the development of optimal transport theory and explains the motivation behind our study. We first recall the evolution from the classical two-marginal formulation to the multimarginal setting, then discuss the importance of the dual formulation and the structural role of canonical dual potentials. We conclude by outlining the main contributions of the paper and its organization.

\subsection{From Monge to Kantorovich, and to several marginals}

Optimal transport (OT) has matured, over more than two centuries, from a
geometric framework to modern analysis,
probability, and machine learning. Its origin is Gaspard Monge's 1781
memoir~\cite{monge1781memoire}, in which the problem of displacing one mass
distribution onto another at least total cost was first posed. Monge's formulation was geometrically appealing, but difficult to study with the analytical tools available at the time. The admissible objects are transport \emph{maps}. The prescribed image-measure constraint is nonlinear, so the resulting variational problem is generally nonconvex and may have no solution.
In~\cite{article}, Kantorovich replaced transport maps by \emph{couplings}, that is, probability measures on the product space with prescribed marginals. This reformulated the problem as an infinite-dimensional linear program. This relaxation is not merely a technical convenience: it
linearizes the constraint set, convexifies the objective, and leads to a duality theory in which the primal transport cost is matched by a supremum over families of potentials. The
refinements of Sudakov~\cite{Sudakov1979}, Knott--Smith~\cite{knott1984}, and
Rachev--Rüschendorf~\cite{rachev1998}, among many others, brought the
two-marginal theory to its present state of completeness, in which the duality
identity, the existence of optimal plans, and the $c$-cyclical structure of their
supports are thoroughly understood.

Extending the problem from two marginals to several is not straightforward and was developed much later. For $K\ge3$ marginals
$\mu_1,\dots,\mu_K$ on spaces $X_1,\dots,X_K$ and a cost
$c:X_1\times\cdots\times X_K\to\mathbb R$, the MOT problem reads
\[
\mathsf{MOT}_c(\mu_1,\dots,\mu_K)
=\inf_{\pi\in\Pi(\mu_1,\dots,\mu_K)}
\int_{X_1\times\cdots\times X_K} c\,d\pi,
\]
where $\Pi(\mu_1,\dots,\mu_K)$ denotes the couplings with $k$-th marginal $\mu_k$.
The foundational existence and structure theory was established by Gangbo and
\'Swiech~\cite{gangbo1998}, and the subject was subsequently developed by
Heinich~\cite{heinich2002}, Carlier and Ekeland~\cite{carlier2010}, and, in a
sequence of contributions, by Pass~\cite{pass2012,pass2015}. What
distinguishes the multimarginal regime is precisely this geometric richness. The two-marginal problem is replaced by $K$ simultaneous marginal constraints. The support of an optimal plan may have a complex geometric structure~\cite{colombo2015}. More importantly for this paper, the dual potentials become harder to characterize.


\subsection{Why the dual side, and why its structure}

In modern optimal transport, the dual formulation is more than a technical tool. It provides a main framework for studying and understanding optimal transport problems. Three reasons make the structure of optimal dual families, rather than only the equality of primal and dual \emph{values}, a central object of study.


First, in statistical and probabilistic applications, the regularity of optimal
potentials is the quantity that propagates into limit theory. The recent
developments in empirical optimal transport, functional delta methods, and
central limit theorems for transport functionals
\cite{del2019,tameling2019,fournier2015} rely on differentiating the optimal
value with respect to the marginals; the derivative is expressed through the dual
potentials, and such expansions depend on their existence, regularity, and uniqueness up to the natural indeterminacy. A duality theorem
that delivers only the value, without a canonical and well-behaved maximizer, is
insufficient for these purposes.

Second, in learning and computation, the dual potentials, or their parameterized approximations, are the quantities that are optimized. MOT provides a natural way to compare and combine several distributions at the same time. It is now widely used in multi-domain adaptation~\cite{Hui2018,He2019}, multi-source generative modeling~\cite{cao2019,genevay2018,deshpande2018}, and barycentric representation learning~\cite{peyre2019}. In adversarial and gradient-based methods, the dual potentials are the quantities being optimized. In barycenter and multi-distribution averaging problems~\cite{peyre2019}, $c$-conjugation and related operations also play a key role in iterative algorithms. Thus, dual potentials are not only analytical tools: they describe the geometry of the problem through $c$-splitting sets and support numerical methods through $c$-transforms. Their existence in a \emph{canonical} form is therefore important for both interpretation and numerical stability.

Third, and underlying both, is a question of mathematical structure intrinsic to
the multimarginal problem. The two-marginal theory possesses a clean notion of
conjugate potentials: an optimal pair $(\varphi,\psi)$ may be chosen so that each
is the $c$-transform of the other, which simultaneously yields equicontinuity,
canonical representatives, and the link to $c$-cyclical monotonicity. The natural
multimarginal analogue replaces this pairwise conjugacy by a \emph{mutual}
conjugacy, in which each component $f_k$ is recovered as the partial $c$-transform
of $f_k$ against \emph{all} the remaining components, a genuinely $K$-fold
coupling with no reduction to the bivariate case. Whether, and under what
hypotheses, the multimarginal Kantorovich dual admits a maximizer in this
mutually $c$-conjugate class is the structural question this paper addresses.

\subsection{Scope of the present work}
The primal-dual \emph{equality} for marginal problems is, in its sharpest
$\sigma$-additive form, already part of the classical theory: Kellerer's
marginal-duality theorem~\cite{kellerer} delivers the equality of values in
considerable generality, and we neither claim it as new nor attempt to improve
upon its hypotheses. Our purpose is different. We
revisit Kellerer-type duality \emph{in the language of MOT}, with bounded continuous costs on product spaces, and we focus on the structure of optimal dual families. In particular, we show that, under suitable conditions, dual maximizers can be chosen as mutually $c$-conjugate potentials. This structure is closely related to $c$-splitting sets and to the supports of optimal primal plans. We treat the compact and non-compact cases separately, since they require different arguments.

In the \emph{compact} setting, the continuity of the cost and the compactness of
the underlying spaces place the problem within the reach of classical convex
duality. We reformulate the dual problem through two convex functionals on the
Banach space of continuous functions on the product and identify the dual via a
Fenchel--Rockafellar argument. This makes the additive structure of multimarginal dual potentials explicit, extending the separable structure of the two-marginal case. The structural refinement then proceeds as
follows. Mutual $c$-conjugacy forces each potential to inherit the modulus of
continuity of the cost, hence equicontinuity; after a normalization that quotients
out the additive indeterminacy of the dual problem, the freedom to shift the
components by constants of vanishing total sum, equicontinuity upgrades to
uniform boundedness. The Arzelà--Ascoli theorem then extracts a uniformly
convergent maximizing sequence whose limit is admissible, optimal, and again
mutually $c$-conjugate. The conclusion is that the dual maximum is attained, and
may be attained within the canonical class.

In the \emph{non-compact} Polish setting, compactness of the dual potentials on the whole space is no longer available, so a different approach is needed. We first recover compactness on the primal side. Since the marginals are tight, the set of couplings is also tight. Prokhorov's theorem then gives narrow compactness of the coupling set. For bounded continuous costs, this implies the existence of an optimal plan and a finite transport value, without any coercivity assumption on the cost. To prove the duality identity, we use a truncation--tightness argument. We first restrict the problem to compact subsets and apply the compact duality result there. The truncated dual potentials are then normalized and adjusted component by component. They remain uniformly bounded by constants depending only on $|c|_\infty$ and $K$. Finally, we extend them to the whole space by assigning a sufficiently negative constant outside the compact sets. This preserves global admissibility, while the change in the dual value tends to zero as the truncation error vanishes. This truncation--normalization--regularization
strategy adapts to the multimarginal framework the classical two-marginal approach
of Villani~\cite{villani2003}. For \emph{dual attainment} in the non-compact case
we turn to the geometry of optimal supports. On the support of an optimal plan the
dual constraint saturates, which is exactly the notion of a $c$-splitting
set~\cite{griessler2016}; we show that the existence of a single real-valued
splitting family can be upgraded, by a maximality (Zorn) argument along the
componentwise order, to a mutually $c$-conjugate splitting family, that boundedness
of the cost forces any such family to be bounded, and hence that, when the family
is Borel measurable, it furnishes a genuine, bounded, Borel measurable dual
optimizer in the canonical class.

\paragraph{Contributions.}
We summarize the principal contributions.

\begin{itemize}
\item \emph{(Compact duality.)} We recover the multimarginal Kantorovich duality
identity for bounded continuous costs on compact product spaces through a
convex-analytic reformulation and the Fenchel--Rockafellar theorem. (Theorem~\ref{KANTO}).

\item \emph{(Non-compact duality.)} We extend the duality identity to general
non-compact Polish product spaces by a quantitative truncation procedure built on
the tightness of multimarginal plans and the boundedness of the cost, with no
coercivity assumption (Theorem~\ref{KANTO2}).

\item \emph{(Canonical dual structure.)} We establish mutually \(c\)-conjugate representations of optimal dual potentials under the corresponding compact and non-compact hypotheses. In the compact case this follows from
conjugation, equicontinuity, normalization, and an Arzelà--Ascoli compactness
argument (Theorem~\ref{KKK}); in the non-compact case it is obtained through
$c$-splitting sets, the supports of optimal plans, and the construction of
bounded Borel measurable mutually $c$-conjugate representatives
(Proposition~\ref{sp}, Proposition~\ref{bounded-splitting},
Theorem~\ref{noncompact}).
\end{itemize}
Taken together, these results identify the mutually $c$-conjugate potentials as
the canonical representatives of the multimarginal Kantorovich dual whenever the
stated compactness or support hypotheses hold. This structure provides a basis for further work on statistical and probabilistic MOT, including stability with respect to the marginals, differentiability of the transport value, central limit theorems, functional delta methods, and numerical approximation.

\paragraph{Organization.}
Section~\ref{wd} fixes the notation, introduces the primal and dual multimarginal
problems, and establishes the weak duality inequality. Section~\ref{secdual}
derives the compact duality identity via the Fenchel--Rockafellar theorem.
Section~\ref{secdual2} extends the identity to non-compact Polish spaces through
tightness and truncation. Section~\ref{secex} studies the structure of optimal
dual potentials, proving the existence of mutually $c$-conjugate maximizers in the
compact case and, via $c$-splitting sets, dual attainment in the non-compact case.
Section~\ref{sec:related} situates the contribution within the literature, and the
Appendix~\ref{secapp} collects the proofs and the supporting results.







\section{Preliminaries and weak duality}\label{wd} 
Let \(K \ge 3\), and let \(X_1,\ldots,X_K\) be topological spaces endowed with their Borel \(\sigma\)-algebras. Let \(\mu_k \in \mathcal{P}(X_k)\) for \(k=1,\dots,K\), where \(\mathcal{P}(X_k)\) denotes the set of probability measures on \(X_k\). For each $k=1,\dots,K$, we denote by $\mathrm{pr}_k: X_1 \times \cdots \times X_K \to X_k$ the canonical projection onto the $k$-th coordinate. We define the set of multimarginal couplings by
\begin{align*}
\Pi(\mu_1,\dots,\mu_K)
= \bigl\{\, \pi \in \mathcal{P}(X_1 \times \cdots \times X_K)
\;|\; (\mathrm{pr}_k)_{\#}\pi = \mu_k \ \text{for all } k=1,\dots,K \,\bigr\}.
\end{align*}
Let $c : X_1 \times \cdots \times X_K \to \mathbb{R}$ be a bounded measurable cost function.
For a transport plan $\pi$ and a family of functions $\boldsymbol{f} = (f_1,\dots,f_K)$
with $f_k : X_k \to \mathbb{R}$, we define
\begin{align*}
I(\pi)
&= \int_{X_1 \times \cdots \times X_K}
c(x_1,\ldots,x_K)\, d\pi(x_1,\ldots,x_K)  \text{  and   } 
J(\boldsymbol{f})
= \sum_{k=1}^{K} \int_{X_k} f_k\, d\mu_k.
\end{align*}
The MOT problem associated with the cost function $c$
and the marginals $\mu_1,\ldots,\mu_K$ is defined as
\begin{equation*}
\mathsf{MOT}_c(\mu_1,\ldots,\mu_K)
= \inf_{\pi \in \Pi(\mu_1,\ldots,\mu_K)} I(\pi).
\end{equation*}
We define the admissible class


\begin{equation*}\label{Fc}
\begin{aligned}
\mathcal F_c
=\;&
\Bigl\{
\boldsymbol f=(f_1,\ldots,f_K)\ \Big|\ 
\forall k,\ f_k:X_k\to\mathbb R \text{ is measurable and }
f_k\in L_1(d\mu_k),
\\
&\qquad
\text{and }\;
\sum_{k=1}^K f_k(x_k)\le c(x_1,\ldots,x_K),
\ \forall (x_1,\ldots,x_K) \in X_1 \times\ldots\times X_K
\Bigr\}.
\end{aligned}
\end{equation*}
We recall that, for each \(k\),
\[
L_1(d\mu_k)
=
\Big\{
f:X_k\to\mathbb R \ \text{measurable}
\;:\;
\int_{X_k}|f|\,d\mu_k<\infty
\Big\},
\]
up to \(\mu_k\)-almost everywhere equality. We also denote by
\[
\mathcal C_b(X_k)
=
\Big\{
f:X_k\to\mathbb R \ \text{continuous and bounded}
\Big\}
\]
the space of bounded continuous functions on \(X_k\). With this notation, the multimarginal Kantorovich primal problem can be written as
\begin{align*}
(\mathsf{MOT}\!-\!\mathsf{KP}):
\inf\bigl\{I(\pi):\;
\pi\in\Pi(\mu_1,\dots,\mu_K)\bigr\},
\end{align*}
and its associated dual problem is given by
\begin{align*}
(\mathsf{MOT}\!-\!\mathsf{DP}):
\sup\Bigl\{
J(\boldsymbol{f}):\;
\boldsymbol{f}\in\mathcal F_c
\Bigr\}.
\end{align*}
With these definitions in place, we now state the weak duality inequality, which holds without any compactness assumption.


\begin{lemma}\label{easyy}
Under the above assumptions on the topological spaces \((X_k)_{k=1}^K\)
endowed with their Borel \(\sigma\)-algebras, the marginals
\((\mu_k)_{k=1}^K\), and the bounded measurable cost function
\(c:X_1\times\cdots\times X_K\to\mathbb R\), the following weak duality relation holds:
\begin{align*}
\sup_{\boldsymbol f\in\mathcal F_c \cap \prod_{k=1}^K \mathcal C_b(X_k)}
J(\boldsymbol f)
\;\le\;
\sup_{\boldsymbol f\in\mathcal F_c}
J(\boldsymbol f)
\;\le\;
\inf_{\pi\in\Pi(\mu_1,\dots,\mu_K)} I(\pi).
\end{align*}
\end{lemma}





























\section{Duality: compact case}\label{secdual}

We first treat the compact case, where the continuity of the cost and the compactness of the underlying spaces allow one to work within the classical framework of convex duality. This setting provides the basic dual formulation of the multimarginal Kantorovich problem and serves as the starting point for the subsequent extensions. The compact duality result is the following.


\begin{theorem}\label{KANTO}
Let $X_1,\ldots,X_K$ be compact metric spaces, and let $\mu_k\in\mathcal{P}(X_k)$.
Let $c : X_1 \times \cdots \times X_K \to \mathbb{R}_+$ be continuous. Then the duality formula holds:
\begin{align*}
\mathsf{MOT}_c(\mu_1,\ldots,\mu_K)
=\inf_{\pi\in\Pi(\mu_1,\dots,\mu_K)} I(\pi)
\;=\;
\sup_{\boldsymbol{f}\in \mathcal{F}_c} J(\boldsymbol{f}).
\end{align*}
\end{theorem}

To prove Theorem~\ref{KANTO}, we reformulate the dual problem using two convex 
functionals defined on the Banach space $E = \mathcal{C}(X_1 \times \cdots \times X_K)$. We now introduce these functionals in two lemmas that gather the properties
required to apply the Fenchel--Rockafellar duality theorem (see Theorem~\ref{fenchell}). 


\begin{lemma}\label{Theta}
The functional $\Theta : E \longrightarrow \mathbb{R}\cup\{+\infty\}$ defined by
\begin{align*}
\Theta(u)=
\begin{cases}
0 & \text{if } u(x_1,\ldots,x_K)\ge -\,c(x_1,\ldots,x_K),\ \\[1mm]
+\infty & \text{otherwise,}
\end{cases}
\end{align*}
is convex and continuous at $u\equiv 1$.
\end{lemma}

\begin{lemma}\label{Xi}
The functional $\Xi:E\to\mathbb{R}\cup\{+\infty\}$ defined by
\begin{align*}
\qquad
\Xi(u)=
\begin{cases}
\displaystyle \sum_{k=1}^{K} \int_{X_k}f_k\,d\mu_k
& \text{if } u(x_1,\ldots,x_K)=\sum_{k=1}^{K} f_k(x_k),\\[2mm]
+\infty & \text{otherwise,}
\end{cases}
\end{align*}
is well defined and convex. Moreover, $\Xi(1)<+\infty$.
\end{lemma}





\section{Duality: non-compact case}\label{secdual2}
In the non-compact case, we must recover a compactness property for the set 
$\Pi(\mu_1,\dots,\mu_K)$. The next two lemmas establish the tightness of multimarginal transference plans and show that this tightness implies compactness under the weak topology. 

\begin{lemma}\label{tight}
Let $X_1,\dots,X_K$ be Polish spaces.  
For each $k\in\{1,\dots,K\}$, let $P_k \subset \mathcal P(X_k)$ be a tight subset of $\mathcal P(X_k)$.  
Then the set $\Pi(P_1,\dots,P_K)$ of all transference plans on $X_1 \times \cdots \times X_K$ whose $k$-th marginal lies in $P_k$ for all $k$, is itself tight in $\mathcal P(X_1 \times \cdots \times X_K)$.
\end{lemma}
\begin{lemma}\label{compact}
Let $X_1,\dots,X_K$ be Polish spaces, and let $\mu_k \in \mathcal P(X_k)$ be fixed probability measures for $k=1,\dots,K$.    
Then $\Pi(\mu_1,\dots,\mu_K)$ is compact in $\mathcal P(X_1 \times \cdots \times X_K)$.
\end{lemma}
\medskip

We now state the non-compact duality theorem. The main point is that, although the spaces \(X_k\) are not compact, the boundedness of the cost and the tightness of the marginals allow one to localize the problem on compact subsets and then pass back to the full space.

\begin{theorem}\label{KANTO2}

Let $X_1,\dots,X_K$ be Polish spaces with $K\geq 3$, and let $\mu_k \in \mathcal P(X_k)$ for $k=1,\dots,K$.  
Let $c : X_1 \times \cdots \times X_K \longrightarrow \mathbb{R}$ be a bounded and continuous cost function. Then the following duality formula holds:
\begin{align*}
\mathsf{MOT}_c(\mu_1,\ldots,\mu_K)
= \inf_{\pi \in \Pi(\mu_1,\dots,\mu_K)} I(\pi)
\;=\;
\sup_{\boldsymbol{f} \in \mathcal F_c}J(\boldsymbol{f}).
\end{align*}
\end{theorem}



\begin{remark}
\item 
Although Theorem~\ref{KANTO} is stated under the standing assumption that the cost
function satisfies \(c\ge 0\), this restriction is inessential on compact product
spaces. In fact, the duality identity remains valid for any bounded and continuous
cost \(c: X_1\times\cdots\times X_K\to\mathbb{R}\), by a simple invariance argument
under additive constants. 
Let \(M\ge -\inf c\) and define the shifted cost \(\widetilde c = c+M\), which is
nonnegative on \(X_1\times\cdots\times X_K\). Since every
\(\pi\in\Pi(\mu_1,\ldots,\mu_K)\) is a probability measure, one has
\begin{align*}
\int \widetilde c\,d\pi
=
\int c\,d\pi + M,
\end{align*}
and therefore
\begin{align*}
\inf_{\pi\in\Pi(\mu_1,\ldots,\mu_K)} \int \widetilde c\,d\pi
=
\inf_{\pi\in\Pi(\mu_1,\ldots,\mu_K)} \int c\,d\pi + M.
\end{align*}
The same translation property holds for the dual problem. Indeed, if
\((f_1,\ldots,f_K)\in\mathcal{F}_c\), then
\((f_1+M,f_2,\ldots,f_K)\in\mathcal{F}_{\widetilde c}\), and since \(\mu_1\) is a
probability measure,
\begin{align*}
J(f_1+M,f_2,\ldots,f_K)
=
J(f_1,\ldots,f_K)+M.
\end{align*}
Conversely, any admissible family
\((g_1,\ldots,g_K)\in\mathcal{F}_{\widetilde c}\) yields
\((g_1-M,g_2,\ldots,g_K)\in\mathcal{F}_c\) with
\begin{align*}
J(g_1-M,g_2,\ldots,g_K)
=
J(g_1,\ldots,g_K)-M.
\end{align*}
Consequently,
\begin{align*}
\sup_{\boldsymbol f\in\mathcal{F}_{\widetilde c}} J(\boldsymbol f)
=
\sup_{\boldsymbol f\in\mathcal{F}_c} J(\boldsymbol f)+M.
\end{align*}
Applying Theorem~\ref{KANTO} to the nonnegative cost \(\widetilde c\) therefore yields,
upon subtracting \(M\), the duality identity for the original bounded continuous cost
\(c\). 
This observation justifies the use of Theorem~\ref{KANTO} on compact truncations of a bounded real-valued cost in the proof of Theorem~\ref{KANTO2}.
\end{remark}









\paragraph{Proof sketch of Theorem~\ref{KANTO2}.}

In the non-compact setting, the main difficulty comes from the lack of compactness of the spaces $X_k$.
 The argument therefore proceeds by a quantitative
localization of the problem onto compact subsets, followed by a stabilization of dual
potentials that allows passage to the limit.

Let $X_k$ be Polish spaces and let $c\in\mathcal C_b(X_1\times\cdots\times X_K)$. Since the
product space is Polish, the set of couplings $\Pi(\mu_1,\dots,\mu_K)$ is narrowly compact by
Prokhorov’s theorem, and the Kantorovich functional $\pi\mapsto\int c\,d\pi$ is narrowly
continuous. Consequently, the primal infimum is finite and is attained by some
$\pi_*\in\Pi(\mu_1,\dots,\mu_K)$. Fix $\varepsilon>0$. By tightness of the marginal measures, there exist compact sets
$X_k^0\subset X_k$ such that $\mu_k(X_k\setminus X_k^0)\le\varepsilon$ for all $k$. Setting
$X^0=X_1^0\times\cdots\times X_K^0$, one obtains the quantitative bound
\(
\pi_*(X\setminus X^0)\le K\varepsilon.
\)
Restricting $\pi_*$ to $X^0$ and renormalizing yields a coupling $\pi_*^0$ with marginals
$\mu_k^0$ supported on $X_k^0$. Since the cost is bounded, the discrepancy between the full
primal value and the truncated primal value is controlled linearly in $\varepsilon$. Hence, the
original problem is approximated, to arbitrary precision, by the compact Kantorovich problem
posed on $X^0$.

On the compact product $X^0$, Kantorovich duality holds by Theorem~\ref{KANTO}. In particular,
for every $\varepsilon>0$, there exist admissible truncated dual potentials
$(\widetilde f_1^0,\dots,\widetilde f_K^0)$ whose dual value is within $\varepsilon$ of the
compact primal optimum. Exploiting the invariance of the admissibility constraint under the
addition of constants with vanishing total sum, we normalize these potentials at a suitable
reference point in $X^0$. This normalization yields uniform upper bounds on each
$\widetilde f_k^0$ over $X_k^0$, with constants depending only on $\|c\|_\infty$ and $K$.

To obtain uniformly bounded representatives while preserving admissibility, we improve the
truncated potentials one by one. More precisely, starting from
\((\widetilde f_1^0,\dots,\widetilde f_K^0)\), we successively replace each
potential by the maximal function compatible with the admissibility inequality
and with the previously constructed potentials. At each step, the admissibility
inequality on \(X^0\) is preserved, and the truncated dual value does not
decrease. The resulting improved potentials
\((\bar f_1^0,\dots,\bar f_K^0)\) are therefore still admissible on \(X^0\),
satisfy $
J^0(\bar f_1^0,\dots,\bar f_K^0)
\ge
J^0(\widetilde f_1^0,\dots,\widetilde f_K^0)$,
and enjoy uniform two-sided bounds depending only on \(K\) and
\(\|c\|_\infty\).

These bounds allow us to extend the improved potentials to the full spaces
while preserving global admissibility. Namely, we keep
\(\bar f_k^0\) on \(X_k^0\), and we define the extension to be a sufficiently
large negative constant outside \(X_k^0\). With this choice, if all coordinates
belong to \(X^0\), admissibility follows from the compact construction; while if
at least one coordinate lies outside its compact set, the large negative value
forces the inequality $
\sum_{k=1}^K f_k(x_k)\le c(x_1,\ldots,x_K)$
on the full product space \(X_1\times\cdots\times X_K\). Thus one obtains a
globally admissible family \((f_1,\dots,f_K)\in\mathcal F_c\).

Since the extended potentials are uniformly bounded, the difference between
their full dual value and the truncated dual value is controlled by the mass of
\(X\setminus X^0\), and hence is of order \(\varepsilon\). Consequently, for
every \(\varepsilon>0\), one constructs a globally admissible family of dual
potentials whose dual value exceeds the primal infimum minus a quantity that
vanishes with \(\varepsilon\). Letting \(\varepsilon\to0\) yields
\[
\sup_{\boldsymbol f\in\mathcal F_c} J(\boldsymbol f)
\;\ge\;
\inf_{\pi\in\Pi(\mu_1,\dots,\mu_K)} I(\pi).
\]
Combined with the converse inequality established in Lemma~\ref{easyy}, this
proves the Kantorovich duality formula in the non-compact Polish setting. This
truncation--normalization--regularization strategy follows the classical
approach introduced by Villani~\cite{villani2003} in the two-marginal case and
extends it to the general multimarginal framework.
















\section{Optimal dual potentials via \texorpdfstring{$c$}{}-conjugacy}\label{secex}


This section refines the duality results obtained in Theorems~\ref{KANTO} and~\ref{KANTO2} by showing that the dual problem admits, under suitable assumptions, optimal potentials with an additional structural property: mutual \(c\)-conjugacy. Such potentials provide canonical representatives of admissible dual families, in the sense that each component is recovered from the others through the corresponding partial \(c\)-transform. We first consider the compact case in which mutual \(c\)-conjugacy yields uniform equicontinuity of the dual potentials. After a normalization which removes the natural additive indeterminacy of the dual problem, one obtains uniform boundedness. The Arzelà--Ascoli theorem then gives compactness of maximizing sequences and leads to the existence of an optimal dual family in the class of mutually \(c\)-conjugate potentials. For the non-compact case,  uniform compactness of potentials on the whole space is no longer available; the proof's argument is based on the structure of optimal primal plans and on \(c\)-splitting sets. This allows us to construct bounded mutually \(c\)-conjugate splitting families on the support of an optimal plan, and hence to recover dual maximizers in the non-compact setting.





\subsection{Compact case}

We begin by recalling the multimarginal analogue of the \(c\)-transform. This definition is inspired by the classical two-marginal case, but here each component is transformed with respect to all the other components of the same family.

\begin{definition}[mutually \(c\)-conjugate]\label{con}
Let $c:X_1\times\cdots\times X_K\longrightarrow \mathbb R$
be a continuous cost function. For a family
\(\boldsymbol f=(f_1,\ldots,f_K)\), with \(f_k:X_k\to\mathbb R\), we define the
partial \(c\)-transform of the \(k\)-th component, with respect to the other
components of the same family, by
\begin{align*}
f_k^{c}(x_k)
=
\inf_{(x_1,\dots,x_{k-1},\,x_{k+1},\dots,x_K)}
\Big\{
c(x_1,\ldots,x_K)
-
\sum_{l\ne k} f_l(x_l)
\Big\},
\qquad \forall x_k \in X_k.
\end{align*}
We say that \(\boldsymbol f=(f_1,\ldots,f_K)\) is a mutually \(c\)-conjugate family if
$f_k=\boldsymbol f_k^{c}$
on $X_k$ for $k=1,\ldots,K$.
\end{definition}


The previous definition gives a pointwise notion of mutual \(c\)-conjugacy. Since the dual problem is formulated on the admissible class \(\mathcal F_c\), we now restrict this notion to admissible dual families.


\begin{definition}\label{Gc}
Using the admissible class \(\mathcal F_c\) defined in \eqref{Fc}, we set
\[
\mathcal G_c
=
\Big\{
\boldsymbol f=(f_1,\ldots,f_K)\in\mathcal F_c:
 f_k=f_k^{c}\ \text{on }X_k
\ \text{for every }k=1,\ldots,K
\Big\}.
\]
Thus \(\mathcal G_c\) is the class of admissible mutually \(c\)-conjugate dual
families.
\end{definition}
\medskip

The next two lemmas provide the compactness tools needed to prove the existence of an optimal dual family in \(\mathcal G_c\). The first lemma shows that partial \(c\)-transforms inherit the modulus of continuity of the cost. The second lemma uses the additive invariance of the dual problem to normalize maximizing sequences without changing either admissibility or the value of the dual functional.



\begin{lemma}\label{equi}
Let \((X_k,d_k)\), \(k=1,\ldots,K\), be compact metric spaces, and let \(c\in C(X_1\times\cdots\times X_K)\). Let \(\omega\) be a modulus of continuity of \(c\). Fix \(k\in\{1,\ldots,K\}\), and let \(f_l:X_l\to\mathbb R\), \(l\ne k\), be such that \[ f_k(x_k) = \inf_{(x_l)_{l\ne k}} \Big\{ c(x_1,\ldots,x_K)-\sum_{l\ne k}f_l(x_l) \Big\} \] is real-valued on \(X_k\). Then 
\[ |f_k(x_k)-f_k(x_k')| \le \omega(d_k(x_k,x_k')), \qquad x_k,x_k'\in X_k. 
\]
\end{lemma}




\begin{lemma}\label{uniform-bounds}
Assume that \(K\ge 3\), that \(X_1,\ldots,X_K\) are compact metric spaces, and that \(c\in C(X_1\times\cdots\times X_K)\), with \(\|c\|_\infty<\infty\).
Let $
(\boldsymbol f^n)_{n\ge1}
=
(f_1^n,\dots,f_K^n)_{n\ge1}
\subset \mathcal G_c$
be a maximizing sequence for $
\sup_{\boldsymbol f\in\mathcal F_c}J(\boldsymbol f)$.
Then there exist constants \(a_1^n,\ldots,a_K^n\in\mathbb R\), with $
\sum_{k=1}^K a_k^n=0$,
such that the translated sequence
\[
\widetilde f_k^n=f_k^n+a_k^n,
\qquad k=1,\ldots,K,
\]
still belongs to \(\mathcal G_c\), is still maximizing, and satisfies the
following uniform bounds: for every \(k=1,\ldots,K\), there exists
\(C_k<\infty\), independent of \(n\), such that
\[
\|\widetilde f_k^n\|_\infty\le C_k
\qquad\text{for all }n.
\]
\end{lemma}


We can now state the compact structural duality result. It shows that, in the compact case, the dual maximum is not only attained in the admissible class \(\mathcal F_c\), but may be attained by a mutually \(c\)-conjugate admissible family. Thus the dual problem can be restricted to the canonical class \(\mathcal G_c\).



\begin{theorem}
\label{KKK}
Let \(X_1,\ldots,X_K\) be compact metric spaces, let
\(\mu_k\in\mathcal P(X_k)\), \(k=1,\ldots,K\), and let $
c:X_1\times\cdots\times X_K\longrightarrow \mathbb R$
be continuous. Then the dual problem \((\mathsf{MOT}\!-\!\mathsf{DP})\)
admits an optimal family $
\boldsymbol f=(f_1,\ldots,f_K)\in\mathcal G_c$.
In particular,
\[
\mathsf{MOT}_c(\mu_1,\ldots,\mu_K)
=
\max_{\boldsymbol f\in\mathcal F_c} J(\boldsymbol f)
=
\max_{\boldsymbol f\in\mathcal G_c} J(\boldsymbol f).
\]
Thus the dual maximum may be attained by an admissible mutually
\(c\)-conjugate family.
\end{theorem}


 



























\subsection{Non-compact case}

We now pass to the non-compact setting. The compactness argument used above cannot be applied directly, since one cannot extract uniformly convergent subsequences of dual potentials on the full spaces \(X_k\). Instead, the argument uses the support of an optimal primal plan. On this support, equality holds in the dual constraint, which naturally leads to the notion of a \(c\)-splitting set. We first recall the definition of a \(c\)-splitting set. Such a set is a subset of the product space on which an admissible dual family saturates the Kantorovich inequality.





\begin{definition}[\(c\)-splitting set \cite{griessler2016}]
Let \(E = X_1 \times \cdots \times X_K\) and  
\(\Gamma \subset E\).
We say that \(\Gamma\) is a \emph{\(c\)-splitting set} if there exist 
\(K\) functions $f_1 : X_1 \to [-\infty,+\infty), 
 \dots,
f_K : X_K \to [-\infty,+\infty)$, such that:
\begin{enumerate}[nosep, leftmargin=*]
\item for all \((x_1,\dots,x_K)\in E\), $\sum_{k=1}^K f_k(x_k) \le c(x_1,\dots,x_K)$.
\item for all \((x_1,\dots,x_K)\in \Gamma\), $\sum_{k=1}^K f_k(x_k) = c(x_1,\dots,x_K)$.
\end{enumerate}
\end{definition}
\medskip

The next proposition shows that, whenever a \(c\)-splitting set admits one real-valued splitting family, one can choose a real-valued splitting family which is maximal with respect to the partial \(c\)-transform operation. In particular, the selected family is mutually \(c\)-conjugate.



\begin{proposition}
\label{sp}
Let \(c:E\to\mathbb R\) be bounded. Let $
\Gamma\subset E=X_1\times\cdots\times X_K$
be a nonempty \(c\)-splitting set. Assume that \(\Gamma\) admits at least one
real-valued splitting family $
(f_1,\ldots,f_K)  \text{; } f_k:X_k\to\mathbb R$.
Then there exists a real-valued splitting family $
(g_1,\ldots,g_K)  \text{; }g_k:X_k\to\mathbb R$,
such that:
\begin{enumerate}
\item for all \((x_1,\ldots,x_K)\in E\) $
\sum_{k=1}^K g_k(x_k)\le c(x_1,\ldots,x_K)$;

\item for all \((x_1,\ldots,x_K)\in\Gamma\) $
\sum_{k=1}^K g_k(x_k)=c(x_1,\ldots,x_K)$;

\item the family is mutually \(c\)-conjugate.
\end{enumerate}
\end{proposition}




The preceding proposition provides the existence of a real-valued mutually
$c$-conjugate splitting family. The next result shows that, when the cost is
bounded, any such family is automatically bounded. Thus, if the family is
Borel measurable, it belongs to the admissible dual class used in the
Kantorovich dual problem.






\begin{proposition}
\label{bounded-splitting}
Let \(X_1,\ldots,X_K\) be arbitrary sets and let
\(
c:X_1\times\cdots\times X_K\to\mathbb R
\)
be bounded. Let
\(\Gamma\subset X_1\times\cdots\times X_K\) be nonempty. Suppose that
\((g_1,\ldots,g_K)\) is a real-valued mutually \(c\)-conjugate splitting family
for \(\Gamma\). Then each \(g_k\) is bounded. In particular, if
\(X_1,\ldots,X_K\) are Polish spaces and the functions \(g_k\) are Borel
measurable, then \((g_1,\ldots,g_K)\) is a bounded Borel measurable mutually
\(c\)-conjugate splitting family.
\end{proposition}

\medskip




With these structural tools at hand, we now turn to dual attainment in the non-compact setting.





\begin{theorem}
\label{noncompact}
Let \(X_1,\ldots,X_K\) be Polish spaces, let
\(\mu_k\in\mathcal P(X_k)\), \(k=1,\ldots,K\), and let $
c:X_1\times\cdots\times X_K\to\mathbb R$
be bounded and continuous. Let $
\pi^*\in\Pi(\mu_1,\ldots,\mu_K)$
be an optimal plan, and set $
\Gamma=\operatorname{supp}(\pi^*)$.
Assume that \(\Gamma\) admits a real-valued Borel measurable mutually
\(c\)-conjugate splitting family. Then the dual problem
\((\mathsf{MOT}\!-\!\mathsf{DP})\) admits an optimizer
\[
\boldsymbol g=(g_1,\ldots,g_K)\in\mathcal G_c .
\]
Moreover,
\[
\mathsf{MOT}_c(\mu_1,\ldots,\mu_K)
=
\min_{\pi\in\Pi(\mu_1,\ldots,\mu_K)}
\int_{X_1\times\cdots\times X_K} c\,d\pi
=
\max_{\boldsymbol f\in\mathcal F_c}J(\boldsymbol f)
=
\max_{\boldsymbol f\in\mathcal G_c}J(\boldsymbol f).
\]
\end{theorem}




\section{Related works}\label{sec:related}

Our results combine three classical approaches: measure-theoretic marginal duality,
existence and regularity of Kantorovich potentials, and the geometry of optimal
supports. We classify each along four axes: additivity class of plans
($\sigma$-additive vs.\ finitely additive)\footnote{A set function $\pi$ is
$\sigma$-additive if, for every sequence of pairwise disjoint measurable sets
$(A_n)_{n\ge1}$, $
\pi\left(\bigcup_{n\ge1}A_n\right)=\sum_{n\ge1}\pi(A_n)$. It is \emph{finitely additive} if the analogous property is required only for
finite collections of pairwise disjoint sets.}, cost class ($\mathcal C_b$, bounded
Borel, or $[0,+\infty]$-valued), number of marginals ($K=2$ vs.\ $K\ge3$), and
\emph{type of dual optimizer} (classical potentials, generalized bidual objects,
or canonical $c$-conjugate families). We position the present paper against
them (Table~\ref{tab:related}).

The primal-dual \emph{equality} for marginal problems is classical
(Kellerer~\cite{kellerer}; Ramachandran--R\"uschendorf under
perfectness\footnote{A probability measure $P$ on a measurable space
$(\Omega,\mathcal A)$ is called \emph{perfect} if, for every
$\mathcal A$-measurable function $f:\Omega\to\mathbb R$, there exists a Borel set
$B\subseteq f(\Omega)$ such that $P(f^{-1}(B))=1$.}~\cite{Ramachandran}), and we do not claim it. Rigo~\cite{Rigo} replaces perfectness
by cost regularity in the two-marginal case; Rigo~\cite{Rigo2} instead
enlarges the admissible class to finitely additive couplings, making attainment
and duality automatic without compactness, tightness, Polishness, or perfectness;
Beiglb\"ock--L\'eonard--Schachermayer~\cite{Beiglb} establish the existence of a dual optimizer for
singular Borel costs via a generalized maximizer in a bidual/projective-limit
space. In every case, the result is either the \emph{optimal value} or an \emph{optimizer} obtained by weakening the original problem (through relaxed plans or bidual potentials).


We move orthogonally on all four axes. We stay $\sigma$-additive, keep
$c\in\mathcal C_b$, work genuinely multimarginally ($K\ge3$, where the partial
$c$-transform of Definition~\ref{con} couples each component to all others), and
produce \emph{canonical} maximizers in the class $\mathcal G_c$ of admissible
mutually $c$-conjugate families. The same cost regularity that elsewhere secures
\emph{values} here drives \emph{structure}: in the compact case the modulus of
$c$ is inherited by every partial $c$-transform (Lemma~\ref{equi}), which after
the zero-sum normalization of Lemma~\ref{uniform-bounds} yields equicontinuous,
uniformly bounded maximizing sequences and hence, by Arzelà--Ascoli, the existence of a maximizer in $\mathcal G_c$ (Theorem~\ref{KKK}); in the non-compact case, we localize via
tightness and Prokhorov compactness (Lemmas~\ref{tight}--\ref{compact},
Theorem~\ref{KANTO2}) and read canonical potentials off the support of an optimal
plan through the $c$-splitting formalism \cite{griessler2016}, upgrading one
splitting family to a Zorn-maximal $c$-conjugate one (Proposition~\ref{sp}),
automatically bounded under bounded $c$ (Proposition~\ref{bounded-splitting}),
yielding a global optimizer in $\mathcal G_c$ (Theorem~\ref{noncompact}). This
support-level structure is exactly where the multimarginal literature
\cite{gangbo1998,pass2012,pass2015,pass2022,colombo2015} enters: we take the
support geometry as given and extract the dual side. Surveys such as
Bogachev--Kolesnikov~\cite{Bogachev} give the global panorama; our statement is
the focused structural complement, and it is precisely the canonical,
perturbation-stable, $c$-transform-computable representative that empirical and
statistical OT \cite{del2019,tameling2019,fournier2015,peyre2019} requires.

\begin{table}[htbp]
\centering
\small
\setlength{\tabcolsep}{5pt}
\renewcommand{\arraystretch}{1.25}
\begin{tabular}{@{}lcccl@{}}
\toprule
\textbf{Work} & \textbf{Additivity} & \textbf{Cost} & \textbf{$K$} &
\textbf{Dual optimizer/ output} \\
\midrule
Kellerer~\cite{kellerer} & $\sigma$-add. & meas. & $\ge2$ &
value only (marginal duality) \\
Ramachandran--R\"usch.~\cite{Ramachandran} & $\sigma$-add. & meas.\ (perfectness) & $2$ &
value only \\
Rigo~\cite{Rigo} & $\sigma$-add. & l.s.c./bdd.\ sections & $2$ &
value; perfectness removed \\
Rigo~\cite{Rigo2} & \emph{finitely add.} & $[0,+\infty]$ Borel & $\ge2$ &
automatic attainment (relaxed plans) \\
Beiglb\"ock et al.~\cite{Beiglb} & $\sigma$-add. & $[0,+\infty]$ Borel & $2$ &
\emph{generalized} bidual maximizer \\
Gangbo--\'Swiech / Pass~\cite{gangbo1998,pass2015} & $\sigma$-add. & smooth/twist & $\ge3$ &
primal support structure \\
\midrule
\textbf{This paper} & $\boldsymbol{\sigma}$\textbf{-add.} & $\boldsymbol{\mathcal C_b}$ &
$\boldsymbol{\ge3}$ & \textbf{canonical }$\boldsymbol{c}$\textbf{-conjugate }$\boldsymbol{\mathcal G_c}$ \\
\bottomrule
\end{tabular}
\caption{Positioning along the four classifying axes. Prior duality is
value-oriented or attains by weakening the problem (finitely additive plans,
bidual potentials); ours stays $\sigma$-additive with bounded continuous cost and
$K\ge3$, and attains \emph{canonical} mutually $c$-conjugate maximizers on both
compact metric and non-compact Polish products.}
\label{tab:related}
\end{table}








\section{Conclusion}\label{conc}








We develop a unified structural treatment of Kantorovich duality for MOT on compact metric and general Polish product spaces with bounded continuous costs. Beyond primal--dual equality, our main contribution is to identify mutually \(c\)-conjugate families as a natural class of dual representatives. In the compact setting, the partial \(c\)-transform inherits the regularity of the cost, and, after normalization, equicontinuity and uniform boundedness yield dual attainment in this class through Arzelà--Ascoli. In the non-compact setting, tightness and localization recover the required compactness for the duality argument, while the geometry of \(c\)-splitting sets provides a route to mutually \(c\)-conjugate dual optimizers under the stated splitting-family assumption. 

These results give a common framework for studying the structure of dual potentials in multimarginal transport and clarify the role played by \(c\)-conjugacy beyond the classical two-marginal setting. The resulting canonical representation is also well suited to further analytical and statistical developments, including stability under perturbations of the marginals and cost, differentiability of the MOT functional, and asymptotic analysis of empirical multimarginal transport.














 
\appendix





% section conclusion (end)


\appendix


\section{Appendix}\label{secapp}
This appendix is divided into three parts in order to separate the auxiliary
arguments from the main exposition and to improve the readability of the paper.
The first part contains additional proofs for the main results and related
auxiliary statements. The second part is devoted to the proofs of the main
duality theorems. The final part collects supporting results that are used
throughout the paper.



\subsection{Additional proofs for the main results}\label{app}

\subsubsection{Proof of Lemma~\ref{easyy}: weak duality}\label{weak}
The inclusion $\mathcal{C}_b(X_k)\subset L_1(\mu_k)$ immediately yields the first inequality.
For the second one, let $\boldsymbol{f}=(f_1,\dots,f_K)\in \mathcal{F}_c$
and $\pi\in \Pi(\mu_1,\dots,\mu_K)$.
By definition of $\Pi$,
\begin{align*}
J(\boldsymbol{f})
=\sum_{k=1}^{K} \int_{X_k} f_k(x_k)\, d\mu_k(x_k) 
=\int_{X_1\times\ldots \times X_K}
\Big( \sum_{k=1}^{K} f_k(x_k) \Big)
\, d\pi(x_1,\ldots,x_K).
\end{align*}
Since $\boldsymbol{f}\in \mathcal{F}_c$, the admissibility condition holds pointwise, namely
\begin{align*}
\sum_{k=1}^K f_k(x_k)\le c(x_1,\dots,x_K)
\qquad\text{for all }(x_1,\dots,x_K)\in X_1\times\cdots\times X_K.
\end{align*}
Therefore,
\begin{equation*}
J(\boldsymbol{f})
=\int_{X_1\times\cdots\times X_K}\sum_{k=1}^K f_k(x_k)\,d\pi(x_1,\dots,x_K)
\le \int_{X_1\times\cdots\times X_K} c\,d\pi
= I(\pi).
\end{equation*}
Thus, for every \(\boldsymbol f\in\mathcal F_c\) and every
\(\pi\in\Pi(\mu_1,\ldots,\mu_K)\), we have
\[
J(\boldsymbol f)\le I(\pi).
\]
Taking the supremum over $\boldsymbol{f}$ on the left-hand side and the infimum over $\pi$ on the right-hand side proves the second inequality.
\hfill\(\Box\)

\subsubsection{Properties of the functionals \texorpdfstring{$\Theta$}{} and \texorpdfstring{$\Xi$}{}}\label{ThetaXi}


\subsubsection*{Proof of Lemma~\ref{Theta}}

Let \(u_1,u_2\in E\) and \(t\in[0,1]\). If \(\Theta(u_1)=\Theta(u_2)=0\), then \(u_1\ge -c\) and \(u_2\ge -c\). Hence \[ t u_1+(1-t)u_2\ge -c, \] and therefore \[ \Theta(tu_1+(1-t)u_2)=0 \le t\Theta(u_1)+(1-t)\Theta(u_2). \] If one of the two quantities \(\Theta(u_1)\), \(\Theta(u_2)\) is equal to \(+\infty\), the convexity inequality is immediate for \(t\in(0,1)\), while the cases \(t=0\) and \(t=1\) are trivial. Thus \(\Theta\) is convex. Since \(c\ge 0\), we have \(1\ge -c\), and hence \(\Theta(1)=0\). Moreover, if \(\|u-1\|_\infty<1\), then \(u>0\ge -c\). Thus \(\Theta(u)=0\) in a neighbourhood of \(1\), which proves that \(\Theta\) is continuous at \(u=1\).
\hfill\(\Box\)

\subsubsection*{Proof of Lemma~\ref{Xi}}
We first prove that \(\Xi\) is well defined. Assume that
\[
\sum_{k=1}^K f_k(x_k)
=
\sum_{k=1}^K \widetilde f_k(x_k),
\qquad
\forall (x_1,\dots,x_K)\in X_1\times\cdots\times X_K,
\]
with \(f_k,\widetilde f_k\in\mathcal C(X_k)\). Set $
d_k=\widetilde f_k-f_k$.
Then $
\sum_{k=1}^K d_k(x_k)=0
\text{ for all }(x_1,\dots,x_K)$.
Fixing all variables except \(x_1\), we obtain that \(d_1\) is constant on
\(X_1\). Repeating the same argument for each coordinate, there exist constants
\(s_1,\dots,s_K\in\mathbb R\) such that
\[
\widetilde f_k-f_k=s_k,
\qquad k=1,\dots,K.
\]
Since \(\sum_{k=1}^K d_k(x_k)=0\), we also have
\[
\sum_{k=1}^K s_k=0.
\]
Therefore, using that each \(\mu_k\) is a probability measure,
\[
\begin{aligned}
\sum_{k=1}^K \int_{X_k} \widetilde f_k\,d\mu_k
&=
\sum_{k=1}^K \int_{X_k} f_k\,d\mu_k
+
\sum_{k=1}^K s_k \mu_k(X_k) \\
&=
\sum_{k=1}^K \int_{X_k} f_k\,d\mu_k
+
\sum_{k=1}^K s_k \\
&=
\sum_{k=1}^K \int_{X_k} f_k\,d\mu_k .
\end{aligned}
\]
Hence \(\Xi\) is well defined.\\
We now prove convexity. Let \(u_1,u_2\in E\) and \(t\in[0,1]\). If
\(\Xi(u_1)=+\infty\) or \(\Xi(u_2)=+\infty\), the convexity inequality is
immediate for \(t\in(0,1)\), while the cases \(t=0\) and \(t=1\) are trivial.
Otherwise, there exist \(h_k,g_k\in\mathcal C(X_k)\) such that
\[
u_1(x_1,\dots,x_K)=\sum_{k=1}^K h_k(x_k),
\qquad
u_2(x_1,\dots,x_K)=\sum_{k=1}^K g_k(x_k).
\]
Then
\[
tu_1+(1-t)u_2
=
\sum_{k=1}^K \bigl(th_k+(1-t)g_k\bigr).
\]
Therefore,
\[
\begin{aligned}
\Xi(tu_1+(1-t)u_2)
&=
\sum_{k=1}^K
\int_{X_k}\bigl(th_k+(1-t)g_k\bigr)\,d\mu_k \\
&=
t\sum_{k=1}^K \int_{X_k} h_k\,d\mu_k
+
(1-t)\sum_{k=1}^K \int_{X_k} g_k\,d\mu_k \\
&=
t\Xi(u_1)+(1-t)\Xi(u_2).
\end{aligned}
\]
Thus \(\Xi\) is convex.\\
Finally, to show that \(\Xi(1)<+\infty\), choose \(f_1\equiv 1\) and
\(f_k\equiv 0\) for every \(k\ge 2\). Then
\[
\Xi(1)=\int_{X_1}1\,d\mu_1=1<+\infty.
\]
This completes the proof.
\hfill\(\Box\)







\subsubsection{Tightness and weak compactness of \texorpdfstring{$\Pi(\mu_1,\ldots,\mu_K)$}{}}\label{tightcompact}

\subsubsection*{Proof of Lemma~\ref{tight}}
Let \(\varepsilon>0\). Since \(P_k\) is tight (see Definition~\ref{tightness}) for each
\(k=1,\ldots,K\), there exists a compact set
\(\mathcal K_{k,\varepsilon}\subset X_k\), independent of the choice of
\(\mu_k\in P_k\), such that
\[
\mu_k(X_k\setminus \mathcal K_{k,\varepsilon})
\le \frac{\varepsilon}{K}
\qquad
\text{for all } \mu_k\in P_k .
\]
Let \(\pi\in \Pi(P_1,\ldots,P_K)\), and denote by \(\pi_k\) its \(k\)-th
marginal. Then \(\pi_k\in P_k\) for every \(k\). Hence
\[
\pi_k(X_k\setminus \mathcal K_{k,\varepsilon})
\le \frac{\varepsilon}{K}.
\]
Therefore,
\[
\begin{aligned}
\pi\Big(
(X_1\times\cdots\times X_K)
\setminus
(\mathcal K_{1,\varepsilon}\times\cdots\times \mathcal K_{K,\varepsilon})
\Big)
&\le
\sum_{k=1}^K
\pi_k(X_k\setminus \mathcal K_{k,\varepsilon})  \\
&\le
\sum_{k=1}^K \frac{\varepsilon}{K}
=
\varepsilon .
\end{aligned}
\]
Since the finite product $
\mathcal K_{1,\varepsilon}\times\cdots\times \mathcal K_{K,\varepsilon}$
is compact in \(X_1\times\cdots\times X_K\), this proves that
\(\Pi(P_1,\ldots,P_K)\) is tight.
\hfill\(\Box\)
   
\subsubsection*{Proof of Lemma~\ref{compact}}
Since each \(X_k\) is Polish, every probability measure on \(X_k\) is tight.
In particular, the singleton \(\{\mu_k\}\) is tight in \(\mathcal P(X_k)\) for
each \(k=1,\ldots,K\).  This fact is known as Ulam’s lemma (see \cite{Krylov2002}, Chap.~1). By Lemma~\ref{tight}, the set
\(\Pi(\mu_1,\ldots,\mu_K)\) is tight in
\(\mathcal P(X_1\times\cdots\times X_K)\).
Moreover, the finite product \(X_1\times\cdots\times X_K\) is Polish.
Therefore, by Prokhorov's theorem, \(\Pi(\mu_1,\ldots,\mu_K)\) is relatively
compact for the topology of weak convergence (see Theorem~\ref{prokhorov}).
It remains to prove that \(\Pi(\mu_1,\ldots,\mu_K)\) is closed. Let
\((\pi_n)_{n\ge1}\subset \Pi(\mu_1,\ldots,\mu_K)\) and assume that
\(\pi_n\) converges weakly to some
\(\pi\in\mathcal P(X_1\times\cdots\times X_K)\) (the convergence is understood in the sense of
Definition~\ref{cv}).  Fix
\(k\in\{1,\ldots,K\}\) and let \(\varphi\in C_b(X_k)\). Since the \(k\)-th
marginal of \(\pi_n\) is \(\mu_k\), we have
\[
\int_{X_1\times\cdots\times X_K}
\varphi(x_k)\,d\pi_n(x_1,\ldots,x_K)
=
\int_{X_k}\varphi\,d\mu_k .
\]
The function
\[
(x_1,\ldots,x_K)\longmapsto \varphi(x_k)
\]
is bounded and continuous on \(X_1\times\cdots\times X_K\). Hence, by weak
convergence,
\[
\int_{X_1\times\cdots\times X_K}
\varphi(x_k)\,d\pi_n
\longrightarrow
\int_{X_1\times\cdots\times X_K}
\varphi(x_k)\,d\pi .
\]
Therefore,
\[
\int_{X_1\times\cdots\times X_K}
\varphi(x_k)\,d\pi
=
\int_{X_k}\varphi\,d\mu_k
\qquad
\text{for all } \varphi\in C_b(X_k).
\]
This means that the \(k\)-th marginal of \(\pi\) is \(\mu_k\). Since \(k\) was
arbitrary, all marginals of \(\pi\) are \(\mu_1,\ldots,\mu_K\), and thus
\[
\pi\in\Pi(\mu_1,\ldots,\mu_K).
\]
Hence \(\Pi(\mu_1,\ldots,\mu_K)\) is closed. Since it is both relatively compact
and closed, it is compact.
\hfill\(\Box\)






\subsubsection{Equicontinuity and normalization of dual potentials}\label{eq_norm}

\subsubsection*{Proof of Lemma~\ref{equi}}
Fix \(x_k,x_k'\in X_k\). For each choice of \((x_l)_{l\ne k}\), define
\[
g_{(x_l)_{l\ne k}}(x_k)
=
c(x_1,\ldots,x_K)-\sum_{l\ne k} f_l(x_l).
\]
The second term does not depend on \(x_k\). Hence
\[
\left|
g_{(x_l)_{l\ne k}}(x_k)
-
g_{(x_l)_{l\ne k}}(x_k')
\right|
=
\left|
c(x_1,\ldots,x_k,\ldots,x_K)
-
c(x_1,\ldots,x_k',\ldots,x_K)
\right|
\le
\omega(d_k(x_k,x_k')).
\]
Therefore,
\[
g_{(x_l)_{l\ne k}}(x_k)
\le
g_{(x_l)_{l\ne k}}(x_k')
+
\omega(d_k(x_k,x_k')).
\]
Taking the infimum over \((x_l)_{l\ne k}\) gives
\[
f_k(x_k)
\le
f_k(x_k')
+
\omega(d_k(x_k,x_k')).
\]
Exchanging \(x_k\) and \(x_k'\) gives the reverse inequality. Hence
\[
|f_k(x_k)-f_k(x_k')|
\le
\omega(d_k(x_k,x_k')).
\]
This proves the result.
\hfill\(\Box\)

\subsubsection*{Proof of Lemma~\ref{uniform-bounds}}
Since \(\boldsymbol f^n\in\mathcal G_c\), each component \(f_k^n\) is a
real-valued partial \(c\)-transform. Therefore, by Lemma~\ref{equi}, for every
\(k=1,\ldots,K\), the family \((f_k^n)_{n\ge1}\) is equicontinuous on \(X_k\),
with a modulus depending only on \(c\). Indeed, all the functions \(f_k^n\) have the same modulus
of continuity \(\omega\), independent of \(n\). Hence, for every
\(x,y\in X_k\),
\[
|f_k^n(x)-f_k^n(y)|
\le \omega(d_k(x,y))
\le \omega(\operatorname{diam}(X_k)).
\]
Since \(X_k\) is compact, \(\operatorname{diam}(X_k)<+\infty\). Therefore,
\[
\sup_{X_k} f_k^n-\inf_{X_k} f_k^n
\le
\omega(\operatorname{diam}(X_k)).
\]
Thus we may take $
D_k=\omega(\operatorname{diam}(X_k))$,
which is independent of \(n\).
For each \(n\), choose \(x_1^{(n)}\in X_1\) such that
\[
f_1^n(x_1^{(n)})=\min_{X_1}f_1^n.
\]
Set
\[
b_1^n=-f_1^n(x_1^{(n)}).
\]
Then
\[
f_1^n(x_1^{(n)})+b_1^n=0.
\]
Fix reference points \(x_k^0\in X_k\), \(k=2,\ldots,K\). Define
\[
\alpha_n
=
\frac{1}{K-1}
\left(
\sum_{k=2}^K f_k^n(x_k^0)-b_1^n
\right),
\]
and set
\[
b_k^n=\alpha_n-f_k^n(x_k^0),
\qquad k=2,\ldots,K.
\]
Then
\[
\sum_{k=1}^K b_k^n
=
b_1^n+\sum_{k=2}^K(\alpha_n-f_k^n(x_k^0))
=
b_1^n+(K-1)\alpha_n-\sum_{k=2}^K f_k^n(x_k^0)
=0.
\]
Define
\[
\widetilde f_k^n=f_k^n+b_k^n,
\qquad k=1,\ldots,K.
\]
Since the added constants have total sum zero, admissibility and the value of
\(J\) are preserved. Indeed, for every \((x_1,\ldots,x_K)\), we have
\[
\sum_{k=1}^K \widetilde f_k^n(x_k)
=
\sum_{k=1}^K f_k^n(x_k)
+
\sum_{k=1}^K b_k^n
=
\sum_{k=1}^K f_k^n(x_k).
\]
Thus the admissibility inequality remains unchanged. Moreover,
\[
J(\widetilde{\boldsymbol f}^{\,n})
=
\sum_{k=1}^K\int_{X_k}(f_k^n+b_k^n)\,d\mu_k
=
J(\boldsymbol f^n)
+
\sum_{k=1}^K b_k^n
=
J(\boldsymbol f^n),
\]
because each \(\mu_k\) is a probability measure and
\(\sum_{k=1}^K b_k^n=0\). Moreover, the mutual \(c\)-conjugacy identities are also
preserved. Indeed, for every \(k\),
\[
\inf_{(x_l)_{l\ne k}}
\Big\{
c(x_1,\ldots,x_K)-\sum_{l\ne k}\widetilde f_l^n(x_l)
\Big\}
=
\inf_{(x_l)_{l\ne k}}
\Big\{
c(x_1,\ldots,x_K)-\sum_{l\ne k}f_l^n(x_l)
\Big\}
-\sum_{l\ne k}b_l^n.
\]
Since \(\boldsymbol f^n\in\mathcal G_c\) and \(\sum_l b_l^n=0\), the last
quantity is $
f_k^n(x_k)+b_k^n
=
\widetilde f_k^n(x_k)$.
Thus \(\widetilde{\boldsymbol f}^{\,n}\in\mathcal G_c\). Also,
\[
J(\widetilde{\boldsymbol f}^{\,n})=J(\boldsymbol f^n),
\]
so the translated sequence is still maximizing.
By construction, $
\widetilde f_1^n(x_1^{(n)})=0$
and $
\widetilde f_k^n(x_k^0)=\alpha_n
\text{ for } k=2,\ldots,K$.
The oscillation bounds are unchanged by translations. Hence
\[
0\le \widetilde f_1^n(x_1)\le D_1,
\qquad x_1\in X_1.
\]
Using admissibility of \(\widetilde{\boldsymbol f}^{\,n}\) at the point
\((x_1^{(n)},x_2^0,\ldots,x_K^0)\), we get
\[
\widetilde f_1^n(x_1^{(n)})
+
\sum_{k=2}^K\widetilde f_k^n(x_k^0)
\le
c(x_1^{(n)},x_2^0,\ldots,x_K^0).
\]
Since \(\widetilde f_1^n(x_1^{(n)})=0\), this gives $
(K-1)\alpha_n\le \|c\|_\infty$.
Thus \((\alpha_n)_n\) is uniformly bounded from above.
It remains to bound \((\alpha_n)_n\) from below. Let
\[
S=\sup_{\boldsymbol f\in\mathcal F_c}J(\boldsymbol f).
\]
The number \(S\) is finite from above, since every admissible
\(\boldsymbol f\) satisfies
\[
J(\boldsymbol f)
\le
\int c\,d(\mu_1\otimes\cdots\otimes\mu_K)
\le
\|c\|_\infty.
\]
Moreover \(S>-\infty\), because the constant family
\[
f_1\equiv-\|c\|_\infty,\qquad f_2=\cdots=f_K\equiv0
\]
is admissible. Since
\((\widetilde{\boldsymbol f}^{\,n})_n\) is a maximizing sequence, after
discarding finitely many terms we may assume
\[
J(\widetilde{\boldsymbol f}^{\,n})\ge S-1.
\]
The finitely many discarded terms can be absorbed into the final constants.
For \(k\ge2\), the oscillation bound gives
\[
\widetilde f_k^n(x_k)
\le
\widetilde f_k^n(x_k^0)+D_k
=
\alpha_n+D_k,
\qquad x_k\in X_k.
\]
Together with \(0\le \widetilde f_1^n\le D_1\), this implies
\[
J(\widetilde{\boldsymbol f}^{\,n})
\le
D_1+\sum_{k=2}^K(\alpha_n+D_k).
\]
Therefore
\[
S-1
\le
D_1+(K-1)\alpha_n+\sum_{k=2}^K D_k.
\]
Hence
\[
\alpha_n
\ge
\frac{S-1-D_1-\sum_{k=2}^K D_k}{K-1}.
\]
Thus \((\alpha_n)_n\) is uniformly bounded.
For \(k\ge2\), the oscillation bound now gives
\[
|\widetilde f_k^n(x_k)|
\le
|\alpha_n|+D_k,
\qquad x_k\in X_k.
\]
The first component satisfies \(0\le\widetilde f_1^n\le D_1\). Therefore, for
every \(k=1,\ldots,K\), there exists \(C_k<\infty\), independent of \(n\), such
that
\[
\|\widetilde f_k^n\|_\infty\le C_k.
\]
This completes the proof.
\hfill\(\Box\)

\subsubsection{Structural properties of splitting families}


\subsubsection*{Proof of Proposition \ref{sp}}
Let \(\mathcal S\) be the set of all real-valued splitting families
\(h=(h_1,\ldots,h_K)\)
for \(\Gamma\) such that
\[
h_k\ge f_k
\qquad\text{on }X_k,\qquad k=1,\ldots,K.
\]
The set \(\mathcal S\) is nonempty, since \((f_1,\ldots,f_K)\in\mathcal S\).
We order \(\mathcal S\) by the componentwise order:
\[
h\le q
\quad\Longleftrightarrow\quad
h_k(x_k)\le q_k(x_k)
\quad\text{for all }x_k\in X_k,\ k=1,\ldots,K.
\]
Let \(\mathcal C\subset\mathcal S\) be a totally ordered chain. For each
\(k=1,\ldots,K\), define
\[
h_k^*(x_k)=\sup_{h\in\mathcal C}h_k(x_k).
\]
We first prove that each \(h_k^*\) is real-valued. Clearly \(h^*\ge f\), so
each \(h_k^*\) is finite from below because \(f_k\) is real-valued. We now prove
that each \(h_k^*\) is finite from above. Fix a point
\[
\bar x=(\bar x_1,\ldots,\bar x_K)\in\Gamma .
\]
Let \(h\in\mathcal C\). Since \(h\) is a splitting family, for every
\(x_k\in X_k\),
\[
h_k(x_k)+\sum_{l\ne k}h_l(\bar x_l)
\le
c(\bar x_1,\ldots,\bar x_{k-1},x_k,\bar x_{k+1},\ldots,\bar x_K).
\]
Since \(h_l\ge f_l\), we have
\[
\sum_{l\ne k}h_l(\bar x_l)
\ge
\sum_{l\ne k}f_l(\bar x_l).
\]
Therefore
\[
h_k(x_k)
\le
c(\bar x_1,\ldots,\bar x_{k-1},x_k,\bar x_{k+1},\ldots,\bar x_K)
-
\sum_{l\ne k}f_l(\bar x_l).
\]
Since \(c\) is bounded, this gives
\[
h_k(x_k)
\le
\|c\|_\infty-\sum_{l\ne k}f_l(\bar x_l).
\]
Taking the supremum over \(h\in\mathcal C\), we obtain
\[
h_k^*(x_k)
\le
\|c\|_\infty-\sum_{l\ne k}f_l(\bar x_l)<+\infty.
\]
Thus each \(h_k^*\) is real-valued.
We claim that \(h^*=(h_1^*,\ldots,h_K^*)\) is again a splitting family for
\(\Gamma\). Fix \((x_1,\ldots,x_K)\in E\). Since \(\mathcal C\) is totally
ordered, for any finite number of elements of \(\mathcal C\), one of them
dominates all the others, hence
\[
\sum_{k=1}^K h_k^*(x_k)
=
\sup_{h\in\mathcal C}\sum_{k=1}^K h_k(x_k).
\]
Indeed, by definition of \(h_k^*\), for every \(h\in\mathcal C\) and every
\(k=1,\ldots,K\), we have
\[
h_k(x_k)\le h_k^*(x_k).
\]
Therefore,
\[
\sum_{k=1}^K h_k(x_k) \le \sum_{k=1}^K h_k^*(x_k).
\]
Taking the supremum over \(h\in\mathcal C\), we obtain
\[
\sup_{h\in\mathcal C}\sum_{k=1}^K h_k(x_k)
\le
\sum_{k=1}^K h_k^*(x_k).
\]
Conversely, let \(\varepsilon>0\). For each \(k=1,\ldots,K\), choose
\(h^k\in\mathcal C\) such that
\[
h_k^*(x_k)\le h_k^k(x_k)+\varepsilon.
\]
Since \(\mathcal C\) is totally ordered, one of the finitely many elements
\(h^1,\ldots,h^K\) dominates all the others; denote it by \(h^0\). Then
\[
\sum_{k=1}^K h_k^*(x_k)
\le
\sum_{k=1}^K h_k^0(x_k)+K\varepsilon
\le
\sup_{h\in\mathcal C}\sum_{k=1}^K h_k(x_k)+K\varepsilon.
\]
Letting \(\varepsilon\downarrow0\) gives the desired equality.
Since every \(h\in\mathcal C\) is splitting, we have
\[
\sum_{k=1}^K h_k(x_k)\le c(x_1,\ldots,x_K).
\]
Taking the supremum gives
\[
\sum_{k=1}^K h_k^*(x_k)\le c(x_1,\ldots,x_K).
\]
Thus the splitting inequality holds.
Now let \((x_1,\ldots,x_K)\in\Gamma\). For every \(h\in\mathcal C\),
\[
\sum_{k=1}^K h_k(x_k)=c(x_1,\ldots,x_K).
\]
Since \(h_k^*\ge h_k\), we get
\[
\sum_{k=1}^K h_k^*(x_k)\ge c(x_1,\ldots,x_K).
\]
Together with the splitting inequality, this gives
\[
\sum_{k=1}^K h_k^*(x_k)=c(x_1,\ldots,x_K)
\qquad\text{on }\Gamma.
\]
Thus \(h^*\) is a splitting family for \(\Gamma\). Since \(h^*\ge f\), we have
\(h^*\in\mathcal S\). Therefore every chain has an upper bound in
\(\mathcal S\). By Zorn's lemma (see Lemma~\ref{Zorn}), \(\mathcal S\) has a
maximal element. Denote it by $
g=(g_1,\ldots,g_K)$.
We now prove that \(g\) is mutually \(c\)-conjugate. Since \(g\) is admissible,
for every \(k\) and every \(x_k\in X_k\),
\[
g_k(x_k)
\le
\inf_{(x_l)_{l\ne k}}
\Big\{
c(x_1,\ldots,x_K)-\sum_{l\ne k}g_l(x_l)
\Big\}.
\]
Define
\[
\widehat g_k(x_k)
=
\inf_{(x_l)_{l\ne k}}
\Big\{
c(x_1,\ldots,x_K)-\sum_{l\ne k}g_l(x_l)
\Big\},
\]
and keep the other components unchanged:
\[
\widehat g_l=g_l,\qquad l\ne k.
\]
By definition of the infimum, the family $
(\widehat g_1,\ldots,\widehat g_K)$
is admissible. Moreover,
\[
\widehat g_k\ge g_k,
\qquad
\widehat g_l=g_l\quad(l\ne k).
\]
We check that the equality on \(\Gamma\) is preserved. Let
\((x_1,\ldots,x_K)\in\Gamma\). Since \(g\) is splitting on \(\Gamma\),
\[
g_k(x_k)+\sum_{l\ne k}g_l(x_l)=c(x_1,\ldots,x_K).
\]
Since \(\widehat g_k\ge g_k\), we obtain
\[
\widehat g_k(x_k)+\sum_{l\ne k}g_l(x_l)\ge c(x_1,\ldots,x_K).
\]
But admissibility of \((\widehat g_1,\ldots,\widehat g_K)\) gives the reverse
inequality. Hence equality holds on \(\Gamma\). Therefore
\[
(\widehat g_1,\ldots,\widehat g_K)\in\mathcal S.
\]
It dominates \(g\). By maximality of \(g\), we must have $
\widehat g_k=g_k$.
Since \(k\) was arbitrary, this holds for every \(k=1,\ldots,K\). Hence
\[
g_k(x_k)
=
\inf_{(x_l)_{l\ne k}}
\Big\{
c(x_1,\ldots,x_K)-\sum_{l\ne k}g_l(x_l)
\Big\},
\qquad k=1,\ldots,K.
\]
Thus \(g\) is mutually \(c\)-conjugate.
\hfill\(\Box\)



\subsubsection*{Proof of Proposition \ref{bounded-splitting}}
Fix a point $
\bar x=(\bar x_1,\ldots,\bar x_K)\in \Gamma$ .
Since \((g_1,\ldots,g_K)\) is real-valued, all the numbers
\(g_l(\bar x_l)\) are finite.
Fix \(k\in\{1,\ldots,K\}\). By mutual \(c\)-conjugacy, for every
\(x_k\in X_k\), we have
\[
g_k(x_k)
=
\inf_{(x_l)_{l\ne k}\in \prod_{l\ne k}X_l}
\Big\{
c(x_1,\ldots,x_K)-\sum_{l\ne k}g_l(x_l)
\Big\}.
\]
Evaluating the infimum at $
(x_l)_{l\ne k}=(\bar x_l)_{l\ne k}$,
we obtain
\[
g_k(x_k)
\le
c(\bar x_1,\ldots,\bar x_{k-1},x_k,\bar x_{k+1},\ldots,\bar x_K)
-
\sum_{l\ne k}g_l(\bar x_l).
\]
Since \(c\) is bounded, this gives a uniform upper bound for \(g_k\).
Hence there exists \(B_k<+\infty\) such that
\[
g_k(x_k)\le B_k,
\qquad x_k\in X_k .
\]
Now fix again \(k\in\{1,\ldots,K\}\). For every \(l\ne k\), the preceding
argument gives
\[
g_l(x_l)\le B_l,
\qquad x_l\in X_l .
\]
Therefore, for every \((x_l)_{l\ne k}\in \prod_{l\ne k}X_l\), we have
\[
c(x_1,\ldots,x_K)-\sum_{l\ne k}g_l(x_l)
\ge
-\|c\|_\infty-\sum_{l\ne k}B_l .
\]
Taking the infimum over \((x_l)_{l\ne k}\in \prod_{l\ne k}X_l\), and using
the mutual \(c\)-conjugacy identity, we obtain
\[
g_k(x_k)
\ge
-\|c\|_\infty-\sum_{l\ne k}B_l .
\]
Thus \(g_k\) is bounded from below. Since \(g_k\) is also bounded from above,
\(g_k\) is bounded.
Since \(k\) was arbitrary, each \(g_k\) is bounded. The final statement
follows immediately from the assumed Borel measurability.
\hfill\(\Box\)



\subsection{Proofs of the main duality theorems}\label{proof}


\subsubsection*{Proof of Theorem~\ref{KANTO}}
We recall from the Riesz representation theorem (see Theorem~\ref{RIEZ}) that the dual space \(E^{\ast}\) of $E=\mathcal{C}(X_1\times\ldots\times X_K)$ can be
identified isometrically with the space 
\(\mathcal{M}(X_{1}\times\cdots\times X_{K})\) of finite
signed Radon measures on the product, endowed with the total variation norm. Hence, to each \(L\in E^{\ast}\) corresponds a unique
\(\pi\in \mathcal{M}(X_{1}\times\cdots\times X_{K})\) such that
\begin{align*}
L(u)=\int_{X_{1}\times\cdots\times X_{K}} u\, d\pi.
\end{align*}
Our goal is to compute both sides of the duality relation provided by the Fenchel--Rockafellar theorem. The left-hand side is clearly
\begin{align*}
&\inf\left\{
\sum_{k=1}^{K} \int_{X_k} f_k\,d\mu_k
\;:\;
f_k\in C(X_k),\ 
\sum_{k=1}^{K} f_k(x_k)\ge -c(x_1,\ldots,x_K)
\right\}
\\
&\qquad =
-\sup\left\{
J(\boldsymbol g)
\;:\;
\boldsymbol g\in \mathcal F_c\cap \prod_{k=1}^{K} C(X_k)
\right\}.
\end{align*}

\paragraph{Legendre-Fenchel transform of \(\Theta\).}
For \(\pi\in \mathcal{M}(X_{1}\times\cdots\times X_{K})\),
\begin{align*}
\Theta^{\ast}(\pi)
&=
\sup_{u\in E}
\Big\{
\int u\, d\pi - \Theta(u)
\Big\}
=
\sup_{\substack{u\in E\\ u\ge -c}}
\int u\, d\pi.
\end{align*}
Hence
\begin{align*}
\Theta^{\ast}(-\pi)
=
\sup_{\substack{u\in E\\ u\le c}}
\int u\, d\pi.
\end{align*}
If \(\pi\) is not a nonnegative measure, then there exists
\(v\in \mathcal{C}_{b}(X_{1}\times\cdots\times X_{K})\), \(v\le 0\), such that
\(\int v\, d\pi>0\). Choosing \(u=t v\) and letting \(t\to +\infty\)
gives \(\Theta^{\ast}(-\pi)=+\infty\). If instead \(\pi\ge 0\), the supremum is attained at \(u=c\), and we obtain $\Theta^{\ast}(-\pi)=\int c\, d\pi$. Thus
\begin{align*}
\Theta^{\ast}(-\pi)
=
\begin{cases}
\displaystyle \int c\, d\pi, & \pi\in \mathcal{M}_{+}(X_{1}\times\cdots\times X_{K}),\\[0.3em]
+\infty, & \text{otherwise}.
\end{cases}
\end{align*}
\paragraph{Legendre-Fenchel transform of \(\Xi\).}
For any \(\pi\in \mathcal{M}(X_{1}\times\cdots\times X_{K})\),
\begin{align*}
\Xi^{\ast}(\pi)
=
\sup_{u\in \mathcal{C}_{b}(X_{1}\times\cdots\times X_{K})}
\Big(
\int u\, d\pi - \Xi(u)
\Big).
\end{align*}
Since \(\Xi(u)<\infty\) only when 
\(
u(x_{1},\ldots,x_{K})=\sum_{k=1}^{K} f_{k}(x_{k})
\),
this becomes
\begin{align*}
\Xi^{\ast}(\pi)
=
\sup_{(f_{1},\ldots,f_{K})}
\Bigg(
\int_{X_{1}\times\cdots\times X_{K}}
\sum_{k=1}^{K} f_{k}(x_{k})\, d\pi
-
\sum_{k=1}^{K} \int_{X_{k}} f_{k} \, d\mu_{k}
\Bigg).
\end{align*}
If
\begin{align*}
\int \sum_{k=1}^{K} f_{k}(x_{k})\, d\pi
=
\sum_{k=1}^{K} \int f_{k}\, d\mu_{k}
\quad\text{for all }(f_{1},\ldots,f_{K}),
\end{align*}
then \(\Xi^{\ast}(\pi)=0\). Otherwise, there exists a family \((f_{1},\ldots,f_{K})\) for which the
difference is nonzero. Scaling \(f_{k}\) by \(t\) (or \(-t\)) yields a
quantity that diverges to \(+\infty\), hence \(\Xi^{\ast}(\pi)=+\infty\). Thus
\begin{align*}
\Xi^{\ast}(\pi)
=
\begin{cases}
0, & \pi\in \Pi(\mu_{1},\ldots,\mu_{K}),\\[0.3em]
+\infty, & \text{otherwise.}
\end{cases}
\end{align*}

Both \(\Theta^{\ast}\) and \(\Xi^{\ast}\) are indicator functions:
\begin{align*}
\Theta^{\ast}(-\pi)=
\begin{cases}
\int c\, d\pi, & \pi\ge 0,\\[0.2em]
+\infty, & \text{otherwise},
\end{cases}
\qquad
\Xi^{\ast}(\pi)=
\begin{cases}
0, & \pi\in\Pi(\mu_{1},\ldots,\mu_{K}),\\[0.2em]
+\infty, & \text{otherwise}.
\end{cases}
\end{align*}
Therefore,
\begin{align*}
\max_{\pi\in E^{\ast}}
\big\{
-\Theta^{\ast}(-\pi)-\Xi^{\ast}(\pi)
\big\}
=
\max_{\pi\in\Pi(\mu_{1},\ldots,\mu_{K})}
\left(
-\int c\, d\pi
\right)
=
-\inf_{\pi\in\Pi(\mu_{1},\ldots,\mu_{K})}
\int c\, d\pi.
\end{align*}
Putting everything together and changing signs, we obtain
\begin{align*}
\inf_{\pi\in\Pi(\mu_{1},\ldots,\mu_{K})} I(\pi)
=
\sup_{\boldsymbol{f}\in \mathcal{F}_{c}\cap\prod_{k=1}^{K}\mathcal{C}_{b}(X_{k})}
J(\boldsymbol{f}).
\end{align*}
By Lemma~\ref{easyy}, this yields
\begin{align*}
\inf_{\pi\in\Pi(\mu_1,\ldots,\mu_K)} I(\pi)
=
\sup_{\boldsymbol f\in \mathcal F_c} J(\boldsymbol f),
\end{align*}
which completes the proof of Theorem~\ref{KANTO}.
\hfill\(\Box\)












\subsubsection*{Proof of Theorem~\ref{KANTO2}}

\medskip
\noindent\emph{Step 0: Preliminary bounds and existence of an optimal plan.}
We define
\[
\|c\|_\infty
=
\sup_{(x_1,\dots,x_K)\in X_1\times\cdots\times X_K}
\big|c(x_1,\dots,x_K)\big|.
\]
In particular, for any probability measure $\pi$ on
$X_1\times\cdots\times X_K$ and any measurable set
$A \subset X_1\times\cdots\times X_K$, one has
\[
\left|\int_A c\, d\pi\right|
\le
\|c\|_\infty\,\pi(A).
\]
Our goal is to reduce the general case to the compact case by a careful
truncation procedure.
Since $c$ is bounded, the value of the Kantorovich problem
\[
\inf_{\pi \in \Pi(\mu_1,\dots,\mu_K)} I(\pi)
\]
is finite. Moreover, this infimum is attained. Indeed, let
$(\pi_n)_{n\ge1}\subset \Pi(\mu_1,\dots,\mu_K)$ be a minimizing sequence, that is,
$
I(\pi_n)\longrightarrow
\inf_{\pi\in\Pi(\mu_1,\dots,\mu_K)} I(\pi).
$
By Lemma~\ref{compact}, the set $\Pi(\mu_1,\dots,\mu_K)$ is compact for the weak
topology. Hence, up to extracting a subsequence, there exists
$\pi_*\in\Pi(\mu_1,\dots,\mu_K)$ such that
$
\pi_n \rightharpoonup \pi_* .
$
Since $c\in \mathcal C_b(X_1\times\cdots\times X_K)$, the map
\[
\pi\longmapsto I(\pi)=\int c\,d\pi
\]
is continuous with respect to weak convergence. Therefore,
\[
I(\pi_*)
=
\lim_{n\to\infty} I(\pi_n)
=
\inf_{\pi\in\Pi(\mu_1,\dots,\mu_K)} I(\pi).
\]
Thus, $\pi_*$ is an optimal transport plan. In the sequel, we fix such a plan
$\pi_* \in \Pi(\mu_1,\dots,\mu_K)$.

\medskip
\noindent\emph{Step 1: Compact truncation and comparison of the primal values.}
 Let $\varepsilon>0$ be such that $K\varepsilon<1$. Since each $X_k$ is Polish, the product space $X_1 \times \cdots \times X_K$ is also Polish (see \cite{bourbaki1998}, Page~195). In particular, each $\mu_k$ is tight (Ulam's lemma). Thus, for each 
$k=1,\dots,K$, there exists a compact subset $X_k^0 \subset X_k$ such that $\mu_k(X_k \setminus X_k^0) \le \varepsilon$. Consequently, the optimal plan $\pi_*$ satisfies
\begin{equation}\label{5}
\pi_*\!\left( (X_1 \times \cdots \times X_K) 
\setminus (X_1^0 \times \cdots \times X_K^0) \right)
\le K\varepsilon .
\end{equation}
Define a new probability measure $\pi^0_*$ supported on $X_1^0 \times \cdots \times X_K^0$ by
\begin{equation}\label{uuu}
\pi^0_*
=
\frac{
\mathbf{1}_{\,X_1^0 \times \cdots \times X_K^0}\,
}{
\pi_*(X_1^0 \times \cdots \times X_K^0)}  \pi_*.
\end{equation}
This is well defined since
$\pi_*(X_1^0 \times \cdots \times X_K^0) > 0$. We define the induced marginals of $\pi_*^0$ by $\mu_k^0 = (\mathrm{pr}_k)_\# \pi_*^0
\text{ for all } k=1,\dots,K$. We denote by
\begin{align*}
\Pi^0(\mu_1^0,\dots,\mu_K^0)
=
\Big\{
\pi^0 \in \mathcal P(X_1^0\times\cdots\times X_K^0)
:\; (\mathrm{pr}_k)_\#\pi^0 = \mu_k^0 \ \text{for all } k
\Big\}
\end{align*}
the set of admissible transference plans on the compact product
$X_1^0\times\cdots\times X_K^0$. Define the truncated functional
\begin{align*}
I^0[\pi^0]
=
\int_{X_1^0 \times \cdots \times X_K^0}
c(x_1,\dots,x_K)\, d\pi^0(x_1,\dots,x_K).
\end{align*}
Let $\widetilde{\pi}^0 \in \Pi^0(\mu_1^0,\dots,\mu_K^0)$ satisfy $I^0[\widetilde{\pi}^0]
=
\inf_{\pi^0 \in \Pi^0(\mu_1^0,\dots,\mu_K^0)} I^0[\pi^0]$. From $\widetilde{\pi}^0$, we construct a plan
$\widetilde{\pi} \in \Pi(\mu_1,\dots,\mu_K)$ in a natural way by gluing together 
$\widetilde{\pi}^0$ with $\pi_*$:
\begin{align*}
\widetilde{\pi}
=
\pi_*(X_1^0 \times \cdots \times X_K^0)\, \widetilde{\pi}^0
\;+\;
\mathbf{1}_{(X_1^0 \times \cdots \times X_K^0)^c}\, \pi_* ,
\end{align*}
refer to Remark \ref{rmq} to see why $\widetilde{\pi} \in \Pi(\mu_1,\dots,\mu_K)$. From the definition of $\widetilde{\pi}$, we have
\begin{align*}
I[\widetilde{\pi}]=
\int_{X_1 \times \cdots \times X_K} c \, d\widetilde{\pi} =
\pi_*(X_1^0 \times \cdots \times X_K^0)
\int_{X_1^0 \times \cdots \times X_K^0} c \, d\widetilde{\pi}^0
\;+\;
\int_{(X_1^0 \times \cdots \times X_K^0)^c} c \, d\pi_*.
\end{align*}
that is,
\begin{align*}
I[\widetilde{\pi}]
=
\pi_*(X_1^0 \times \cdots \times X_K^0)\, I^0[\widetilde{\pi}^0]
\;+\;
\int_{(X_1^0 \times \cdots \times X_K^0)^c}
c(x_1,\dots,x_K)\, d\pi_*(x_1,\dots,x_K).
\end{align*}
Since $0 \leq \pi_*(X_1^0 \times \cdots \times X_K^0) \leq 1$, we write
\begin{align*}
I[\widetilde{\pi}]
- I^0[\widetilde{\pi}^0]
&=
\big(\pi_*(X_1^0 \times \cdots \times X_K^0)-1\big)\,
I^0[\widetilde{\pi}^0]
\;+\;
\int_{(X_1^0 \times \cdots \times X_K^0)^c} c\, d\pi_* .
\end{align*}
Taking absolute values yields
\begin{align*}
\big|I[\widetilde{\pi}] - I^0[\widetilde{\pi}^0]\big|
&\le
\big(1-\pi_*(X_1^0 \times \cdots \times X_K^0)\big)\,\big|I^0[\widetilde{\pi}^0]\big|
\;+\;
\Big|\int_{(X_1^0 \times \cdots \times X_K^0)^c} c\, d\pi_*\Big|.
\end{align*}
Since $|\int_A c\,d\pi|\le \|c\|_\infty\,\pi(A)$ for any probability measure $\pi$ and any measurable set $A$, 
we have in particular
\begin{align*}
\big|I^0[\widetilde{\pi}^0]\big| =
\Big|\int_{X_1^0\times\cdots\times X_K^0} c\,d\widetilde{\pi}^0\Big|
\le \|c\|_\infty
\qquad\text{and}\qquad 
\Big|\int_{(X^0)^c} c\, d\pi_*\Big|
\le \|c\|_\infty\,\pi_*((X^0)^c),
\end{align*}
where $X^0=X_1^0\times\cdots\times X_K^0$. Using moreover \eqref{5}, i.e.
$\pi_*((X^0)^c)\le K\varepsilon$, we obtain
\begin{align*}
\big|I[\widetilde{\pi}] - I^0[\widetilde{\pi}^0]\big|
&\le
\|c\|_\infty\,\pi_*((X^0)^c)
+\|c\|_\infty\,\pi_*((X^0)^c)
\le 2K\|c\|_\infty\,\varepsilon.
\end{align*}
In particular, $ I[\widetilde{\pi}]\le I^0[\widetilde{\pi}^0] + 2K\|c\|_\infty\,\varepsilon.$
Thus,
\begin{align*}
\inf_{\pi\in\Pi(\mu_1,\ldots,\mu_K)} I(\pi)
\;\le\;
\inf_{\pi^0\in\Pi^0(\mu_1^0,\ldots,\mu_K^0)} I^0[\pi^0]
\;+\;
2K\|c\|_\infty\,\varepsilon.
\end{align*}

\medskip
\noindent\emph{Step 2: Choice and normalization of a near-optimal truncated dual family.}
For a $K$-family of functions $(f_1^0,\dots,f_K^0)$ with 
$f_k^0 \in L_1(d\mu_k^0)$, we define the truncated dual functional
\begin{align*}
J^0(f_1^0,\dots,f_K^0)
=
\sum_{k=1}^K \int_{X_k^0} f_k^0 \, d\mu_k^0.
\end{align*}
By Theorem \ref{KANTO} (Kantorovich duality on the compact product
$X_1^0 \times \cdots \times X_K^0$), we know that $\inf I^0
=
\sup J^0$. Thus, for any $\varepsilon>0$, there exists an admissible family 
$(\widetilde f_1^0,\ldots,\widetilde f_K^0) \in \prod_{k=1}^K \mathcal C_b(X_k^0)$ such that $J^0(\widetilde{f}_1^0,\dots,\widetilde{f}_K^0)
>
\sup J^0 - \varepsilon$. Our problem is to construct from the truncated family 
$(\widetilde{f}_1^0,\dots,\widetilde{f}_K^0)$ a new family 
$(f_1,\dots,f_K)$ which is effective for the maximization 
of the full dual functional
\begin{align*}
J(f_1,\dots,f_K)
=
\sum_{k=1}^K \int_{X_k} f_k \, d\mu_k.
\end{align*}
It will be useful to ensure that the admissibility inequality
\begin{align*}\label{admiss}
\widetilde{f}_1^0(x_1) + \cdots + \widetilde{f}_K^0(x_K)
\;\le\;
c(x_1,\dots,x_K)
\end{align*}
is valid for all 
\(
(x_1,\dots,x_K) \in X_1^0 \times \cdots \times X_K^0.
\)
Without loss of generality, we assume that $\varepsilon\le 1$. 
Since $c$ is bounded, we have $
-\|c\|_\infty \le c(x_1,\ldots,x_K) \text{ for all }(x_1,\ldots,x_K)\in X_1^0\times\cdots\times X_K^0$.
Consider the constant family $
(f_1^0,\ldots,f_K^0)
=
(-\|c\|_\infty,0,\ldots,0)$.
Then, for every $(x_1,\ldots,x_K)\in X_1^0\times\cdots\times X_K^0$, we have
\[
\sum_{k=1}^K f_k^0(x_k)
=
-\|c\|_\infty
\le
c(x_1,\ldots,x_K).
\]
Thus, this family is admissible for the truncated dual problem. Its dual value is
\[
J^0(f_1^0,\ldots,f_K^0)
=
\int_{X_1^0} -\|c\|_\infty\,d\mu_1^0
+
\sum_{k=2}^K \int_{X_k^0} 0\,d\mu_k^0.
\]
Since $\mu_1^0$ is a probability measure, this gives
\[
J^0(f_1^0,\ldots,f_K^0)
=
-\|c\|_\infty.
\]
Therefore, since $\sup J^0$ is the supremum over all admissible families, it is at least the value of this particular admissible family. Hence
\[
\sup J^0 \ge -\|c\|_\infty.
\]
Consequently,
\begin{align*}
J^0(\widetilde{f}_1^0,\dots,\widetilde{f}_K^0)
\ge \sup J^0 - \varepsilon 
\ge -\|c\|_\infty-\varepsilon 
\ge -\|c\|_\infty-1.
\end{align*}
For any $\pi^0 \in \Pi^0(\mu_1^0,\dots,\mu_K^0)$, we can write
\begin{align*}
J^0(\widetilde{f}_1^0,\dots,\widetilde{f}_K^0)
=\sum_{k=1}^K \int_{X_k^0} \widetilde{f}_k^0\, d\mu_k^0
=\int_{X_1^0\times\cdots\times X_K^0}
\big(\widetilde{f}_1^0(x_1)+\cdots+\widetilde{f}_K^0(x_K)\big)
\, d\pi^0(x_1,\dots,x_K).
\end{align*}
Since this last integral is $\ge -\|c\|_\infty-1$, we deduce that there exists at least one point
$(x_1^0,\dots,x_K^0)\in X_1^0\times\cdots\times X_K^0$ such that
$\widetilde{f}_1^0(x_1^0)+\cdots+\widetilde{f}_K^0(x_K^0)
\ge -\|c\|_\infty-1$. Define $a_k=\widetilde{f}_k^0(x_k^0) \text{ for } k=1,\dots,K$, and let $m=\frac{1}{K}\sum_{k=1}^K a_k$. Since $\widetilde{f}_1^0(x_1^0)+\cdots+\widetilde{f}_K^0(x_K^0)
\ge -\|c\|_\infty-1$, it follows that $m\ge -(\|c\|_\infty+1)/K$.
We now replace $\widetilde{f}_k^0(x_k)$ by $\widetilde{f}_k^0(x_k)+s_k$ where $s_k=m-a_k$. Since
\begin{equation*}
\sum_{k=1}^K s_k=\sum_{k=1}^K (m-a_k)=Km-\sum_{k=1}^K a_k = 0,
\end{equation*}
the admissibility condition is preserved:
\begin{equation*}
\sum_{k=1}^K (\widetilde{f}_k^0(x_k)+s_k)
=
\sum_{k=1}^K \widetilde{f}_k^0(x_k)
+\sum_{k=1}^K s_k
=
\sum_{k=1}^K \widetilde{f}_k^0(x_k)
\le c(x_1,\dots,x_K).
\end{equation*}
Similarly, the dual functional is unchanged:
\begin{equation*}
J^0(\widetilde{f}_1^0+s_1,\dots,\widetilde{f}_K^0+s_K)
=
\sum_{k=1}^K\int_{X_k^0}\widetilde{f}_k^0\, d\mu_k^0
+\sum_{k=1}^K s_k
=
J^0(\widetilde{f}_1^0,\dots,\widetilde{f}_K^0).
\end{equation*}
Finally, evaluating $(\widetilde{f}_1^0,\ldots,\widetilde{f}_K^0) $ at the reference point 
$(x_1^0,\dots,x_K^0)$ gives
\begin{equation*}
\widetilde{f}_k^0(x_k^0)+(m-a_k)
=
m
\qquad\text{for all } k=1,\dots,K.
\end{equation*}
After adding the constants $s_k$, we still denote the new functions by $\widetilde{f}_k^0$. Combined with $m\ge -(\|c\|_\infty+1)/K$, this yields the uniform lower bound
\begin{equation}\label{rr}
\widetilde{f}_k^0(x_k^0)\ge -(\|c\|_\infty+1)/K,\qquad k=1,\dots,K.
\end{equation}
This implies that, for all $x_k\in X_k^0$,
\begin{align*}
\widetilde{f}_k^0(x_k)
&\leq
c(x_1^0,\dots,x_{k-1}^0,\,x_k,\,x_{k+1}^0,\dots,x_K^0)
-
\sum_{l\ne k}\widetilde{f}_l^0(x_l^0)
\\
&\leq
c(x_1^0,\dots,x_{k-1}^0,\,x_k,\,x_{k+1}^0,\dots,x_K^0)
+
\frac{(K-1)(\|c\|_\infty+1)}{K}.
\end{align*}
Therefore,
\begin{align}\label{tild}
\widetilde{f}_k^0(x_k)
\;\le\; 2\|c\|_\infty +1 .
\end{align}

\medskip
\noindent\emph{Step 3: Improvement of the truncated potentials and uniform bounds.}
We now improve the potentials one by one. Set $
(f_1^{(0)},\ldots,f_K^{(0)})
=
(\widetilde f_1^0,\ldots,\widetilde f_K^0)$. Here and throughout this construction, the superscript in parentheses denotes the iteration index. Thus \(f_k^{(j)}\) denotes the \(k\)-th potential after \(j\) improvement steps.
Assume that \((f_1^{(k-1)},\ldots,f_K^{(k-1)})\) has been constructed and is
admissible on \(X^0=X_1^0\times\cdots\times X_K^0\). We define
\[
f_k^{(k)}(x_k)
=
\inf_{(x_l)_{l\ne k}\in \prod_{l\ne k}X_l^0}
\Big\{
c(x_1,\ldots,x_K)
-
\sum_{l\ne k}f_l^{(k-1)}(x_l)
\Big\},
\]
and we set $
f_l^{(k)}=f_l^{(k-1)}
\text{ for }l\ne k$.
By admissibility of \((f_1^{(k-1)},\ldots,f_K^{(k-1)})\), we have $
f_k^{(k-1)}\le f_k^{(k)}
\text{ on }X_k^0$.
Hence
\[
J^0(f_1^{(k)},\ldots,f_K^{(k)})
\ge
J^0(f_1^{(k-1)},\ldots,f_K^{(k-1)}).
\]
Moreover, by the definition of \(f_k^{(k)}\), the new family
\((f_1^{(k)},\ldots,f_K^{(k)})\) is still admissible on \(X^0\). Indeed, for all
\((x_1,\ldots,x_K)\in X^0\),
\[
f_k^{(k)}(x_k)
\le
c(x_1,\ldots,x_K)
-
\sum_{l\ne k}f_l^{(k)}(x_l),
\]
and therefore $
\sum_{l=1}^K f_l^{(k)}(x_l)\le c(x_1,\ldots,x_K)$.
After performing this construction for \(k=1,\ldots,K\), we set
\[
\bar f_k^0=f_k^{(K)},
\qquad k=1,\ldots,K.
\]
Then \((\bar f_1^0,\ldots,\bar f_K^0)\) is admissible on \(X^0\), and $
J^0(\bar f_1^0,\ldots,\bar f_K^0)
\ge
J^0(\widetilde f_1^0,\ldots,\widetilde f_K^0)$.
Moreover, the improved potentials remain uniformly controlled at the reference
points. Indeed, at the initial step, by \eqref{rr}, we have
\[
f_l^{(0)}(x_l^0)
=
\widetilde f_l^0(x_l^0)
\ge
-\frac{\|c\|_\infty+1}{K},
\qquad l=1,\ldots,K.
\]
Now, during the construction, at each step we replace only one potential, say
\(f_k^{(k-1)}\), by \(f_k^{(k)}\), and we have already proved that
\[
f_k^{(k-1)}\le f_k^{(k)}
\qquad\text{on }X_k^0.
\]
All the other potentials are left unchanged. Hence no potential decreases on
the compact set during the construction. Therefore, the lower bound at the
reference points is preserved at every step:
\[
f_l^{(k-1)}(x_l^0)
\ge
-\frac{\|c\|_\infty+1}{K},
\qquad l=1,\ldots,K,\quad k=1,\ldots,K.
\]
Therefore, by evaluating the infimum defining \(f_k^{(k)}\) at
\((x_l)_{l\ne k}=(x_l^0)_{l\ne k}\), we obtain, for every \(x_k\in X_k^0\),
\[
f_k^{(k)}(x_k)
\le
c(x_1^0,\ldots,x_{k-1}^0,x_k,x_{k+1}^0,\ldots,x_K^0)
-
\sum_{l\ne k}f_l^{(k-1)}(x_l^0)
\le
2\|c\|_\infty+1.
\]
Thus each improved potential is uniformly bounded from above.
For the lower bound, we argue by induction. At the initial step, the functions
\(\widetilde f_l^0=f_l^{(0)}\) satisfy the upper bound \eqref{tild}. Moreover,
at each improvement step, the newly defined function also satisfies the same
upper bound. Hence, for every \(k=1,\ldots,K\), all the functions
\(f_l^{(k-1)}\), \(l\ne k\), which appear in the infimum are bounded from above
by \(2\|c\|_\infty+1\). Therefore, for every choice of \((x_l)_{l\ne k}\),
\[
c(x_1,\ldots,x_K)
-
\sum_{l\ne k}f_l^{(k-1)}(x_l)
\ge
-\|c\|_\infty-(K-1)(2\|c\|_\infty+1).
\]
Taking the infimum gives $
f_k^{(k)}(x_k)
\ge
-(2K-1)\|c\|_\infty-(K-1)$.
Consequently, after the \(K\)-th step, there exists a constant \(C>0\),
depending only on \(K\) and \(\|c\|_\infty\), such that
\[
|\bar f_k^0(x_k)|\le C,
\qquad x_k\in X_k^0,\quad k=1,\ldots,K.
\]
\medskip
\noindent\emph{Step 4: Extension to a globally admissible family.}
We now modify the functions outside the compact sets in order to obtain a
globally admissible family. Let $
M=(K-1)C+\|c\|_\infty$.
For every \(k=1,\ldots,K\), define
\[
f_k(x_k)
=
\begin{cases}
\bar f_k^0(x_k), & x_k\in X_k^0,\\
-M, & x_k\in X_k\setminus X_k^0.
\end{cases}
\]
Moreover, the functions \(\bar f_k^0\) are measurable on \(X_k^0\). Indeed,
the initial functions \(\widetilde f_k^0\) are continuous on \(X_k^0\). At each
improvement step, the new potential is obtained as an infimum, over the compact
set \(\prod_{l\ne k}X_l^0\), of a continuous function. Hence the improved
potentials remain continuous, and therefore measurable, on their compact
domains.
Since \(X_k^0\) is compact in the Polish space \(X_k\), it is a Borel subset of
\(X_k\). Hence both \(X_k^0\) and \(X_k\setminus X_k^0\) are Borel sets. On
\(X_k^0\), the function \(f_k\) coincides with the measurable function
\(\bar f_k^0\), while on \(X_k\setminus X_k^0\), it is equal to the constant
\(-M\). Therefore \(f_k\) is measurable on \(X_k\).
Moreover, each \(f_k\) is bounded, so
\(f_k\in L_1(\mu_k)\).
We claim that $
(f_1,\ldots,f_K)\in\mathcal F_c$.
Indeed, let \((x_1,\ldots,x_K)\in X_1\times\cdots\times X_K\). If
\((x_1,\ldots,x_K)\in X_1^0\times\cdots\times X_K^0\), then, since
\((\bar f_1^0,\ldots,\bar f_K^0)\) is admissible on \(X^0\),
\[
\sum_{k=1}^K f_k(x_k)
=
\sum_{k=1}^K \bar f_k^0(x_k)
\le c(x_1,\ldots,x_K).
\]
If, on the other hand, at least one coordinate \(x_j\) does not belong to
\(X_j^0\), then \(f_j(x_j)=-M\), while \(f_k(x_k)\le C\) for every
\(k\ne j\). Hence
\[
\sum_{k=1}^K f_k(x_k)
\le
-M+(K-1)C
=
-\|c\|_\infty
\le
c(x_1,\ldots,x_K).
\]
Thus
\[
\sum_{k=1}^K f_k(x_k)\le c(x_1,\ldots,x_K)
\qquad
\text{for all }(x_1,\ldots,x_K)\in X_1\times\cdots\times X_K.
\]
Therefore \((f_1,\ldots,f_K)\in\mathcal F_c\).

\medskip
\noindent\emph{Step 5: Comparison of the full and truncated dual values.}
Once we have these bounds and the globally admissible family
\((f_1,\ldots,f_K)\), we are almost done. Indeed,
\[
J(f_1,\ldots,f_K)
=
\sum_{k=1}^K \int_{X_k} f_k\,d\mu_k
=
\int_{X_1\times\cdots\times X_K}
\left(\sum_{k=1}^K f_k(x_k)\right)\,d\pi_* .
\]
We now split the product space into the two disjoint parts $
X_1\times\cdots\times X_K
=
X^0\cup (X^0)^c$.
Thus,
\[
J(f_1,\ldots,f_K)
=
\int_{X^0}
\left(\sum_{k=1}^K f_k(x_k)\right)d\pi_*
+
\int_{(X^0)^c}
\left(\sum_{k=1}^K f_k(x_k)\right)d\pi_*.
\]
On \(X^0\), we have \(x_k\in X_k^0\) for every \(k=1,\ldots,K\). By the
definition of the extension, \(f_k=\bar f_k^0\) on \(X_k^0\). Hence, on \(X^0\),
\[
\sum_{k=1}^K f_k(x_k)
=
\sum_{k=1}^K \bar f_k^0(x_k).
\]
Consequently,
\[
\int_{X^0}
\left(\sum_{k=1}^K f_k(x_k)\right)d\pi_*
=
\int_{X^0}
\left(\sum_{k=1}^K \bar f_k^0(x_k)\right)d\pi_*.
\]
By definition of the normalized restriction \(\pi_*^0\), we have
\[
\pi_*^0
=
\frac{\mathbf 1_{X^0}}{\pi_*(X^0)}\,\pi_*.
\]
Equivalently, on \(X^0\),
\[
d\pi_*=\pi_*(X^0)\,d\pi_*^0.
\]
Therefore,
\[
\int_{X^0}
\left(\sum_{k=1}^K \bar f_k^0(x_k)\right)d\pi_*
=
\pi_*(X^0)
\int_{X^0}
\left(\sum_{k=1}^K \bar f_k^0(x_k)\right)d\pi_*^0.
\]
Combining the previous identities, we obtain
\[
J(f_1,\ldots,f_K)
=
\pi_*(X^0)
\int_{X^0}
\left(\sum_{k=1}^K \bar f_k^0(x_k)\right)d\pi_*^0
+
\int_{(X^0)^c}
\left(\sum_{k=1}^K f_k(x_k)\right)d\pi_* .
\]
Moreover,
\[
J^0(\bar f_1^0,\ldots,\bar f_K^0)
=
\int_{X^0}
\left(\sum_{k=1}^K \bar f_k^0(x_k)\right)d\pi_*^0 .
\]
Since \(|\bar f_k^0|\le C\) and \(|f_k|\le M\), and since
\(\pi_*((X^0)^c)\le K\varepsilon\), we obtain
\[
\left|
J(f_1,\ldots,f_K)
-
J^0(\bar f_1^0,\ldots,\bar f_K^0)
\right|
\le
K^2(C+M)\varepsilon .
\]
In particular,
\[
J(f_1,\ldots,f_K)
\ge
J^0(\bar f_1^0,\ldots,\bar f_K^0)
-
K^2(C+M)\varepsilon .
\]
Since
\[
J^0(\bar f_1^0,\ldots,\bar f_K^0)
\ge
J^0(\widetilde f_1^0,\ldots,\widetilde f_K^0)
\ge
\inf I^0-\varepsilon
\]
and
\[
\inf I^0
\ge
\inf I-2K\|c\|_\infty\,\varepsilon,
\]
we deduce that
\[
J(f_1,\ldots,f_K)
\ge
\inf I
-
\Bigl(2K\|c\|_\infty+1+K^2(C+M)\Bigr)\varepsilon .
\]

\medskip
\noindent\emph{Step 6: Passage to the limit and conclusion.}
Since \((f_1,\ldots,f_K)\in\mathcal F_c\), we have $
\sup_{\boldsymbol f\in\mathcal F_c}J(\boldsymbol f)
\ge
J(f_1,\ldots,f_K)$.
Therefore,
\[
\sup_{\boldsymbol f\in\mathcal F_c}J(\boldsymbol f)
\ge
\inf I
-
\Bigl(2K\|c\|_\infty+1+K^2(C+M)\Bigr)\varepsilon .
\]
Letting \(\varepsilon\to0\), we obtain
\[
\sup_{\boldsymbol f\in\mathcal F_c}J(\boldsymbol f)
\ge
\inf_{\pi\in\Pi(\mu_1,\ldots,\mu_K)} I(\pi).
\]
Combining this inequality with the weak duality inequality proved in
Lemma~\ref{easyy}, we conclude that
\[
\inf_{\pi\in\Pi(\mu_1,\ldots,\mu_K)} I(\pi)
=
\sup_{\boldsymbol f\in\mathcal F_c}J(\boldsymbol f).
\]
This proves the Kantorovich duality formula in the non-compact setting.
\hfill\(\Box\)














\begin{remark}\label{rmq}
\noindent
\item To verify that the plan \(\widetilde{\pi}\) defined by $\widetilde{\pi}=\pi_*(X^0)\,\widetilde{\pi}^0+\mathbf{1}_{(X^0)^c}\,\pi_*
\text{ where } X^0=X_1^0\times\cdots\times X_K^0$
belongs to \(\Pi(\mu_1,\dots,\mu_K)\), it is enough to check that each of its marginals coincides with the corresponding $\mu_k$. Fix a measurable set \(A\subset X_k\). Using the definition of
\(\widetilde{\pi}\), we write
\begin{eqnarray*}
(\widetilde{\pi})_k(A)
&=& \widetilde{\pi}(X_1\times\cdots\times A\times\cdots\times X_K)
\\
&=& \pi_*(X^0)\,
     \widetilde{\pi}^0\!\big((X_1\times\cdots\times A\times\cdots\times X_K)\cap X^0\big)
\\
&&
   +\pi_*\!\big((X^0)^c\cap(X_1\times\cdots\times A\times\cdots\times X_K)\big).
\end{eqnarray*}
Since \(\widetilde{\pi}^0\) is supported on \(X^0\) and has \(k\)-th marginal \(\mu_k^0\),
the first term equals
\begin{eqnarray*}
\pi_*(X^0)\,\mu_k^0(A\cap X_k^0)
&=&\pi_*(X^0)\,
\frac{\pi_*\!\big((X_1^0\times\cdots\times(A\cap X_k^0)\times\cdots\times X_K^0)\big)}{\pi_*(X^0)}\\
&=&\pi_*\!\big((X_1\times\cdots\times A\times\cdots\times X_K)\cap X^0\big),
\end{eqnarray*}
where we used \eqref{uuu}. Therefore,
\begin{align*}
(\widetilde{\pi})_k(A)
&=
\pi_*\big((X_1\times\cdots\times A\times\cdots\times X_K)\cap X^0\big)
\\
&\qquad+
\pi_*\big((X^0)^c\cap(X_1\times\cdots\times A\times\cdots\times X_K)\big).
\end{align*}
The two sets in the last line form a disjoint partition of the cylinder
\(X_1\times\cdots\times A\times\cdots\times X_K\); hence
\begin{align*}
(\widetilde{\pi})_k(A)
=\pi_*\big(X_1\times\cdots\times A\times\cdots\times X_K\big)
=\mu_k(A)
\end{align*}
because \(\pi_*\in\Pi(\mu_1,\dots,\mu_K)\). Thus, each marginal of \(\widetilde{\pi}\) is
equal to \(\mu_k\), and we conclude that \(\widetilde{\pi}\in\Pi(\mu_1,\dots,\mu_K)\).

\end{remark}

















\subsubsection*{Proof of Theorem~\ref{KKK}}

\noindent\emph{Finiteness of the dual value.}
Since $c$ is continuous on the compact product $X_1\times\cdots\times X_K$, it is
bounded; set $M=\|c\|_\infty<\infty$ and
$V_c=\sup_{\boldsymbol f\in\mathcal F_c}J(\boldsymbol f)$. If
$\boldsymbol f=(f_1,\dots,f_K)\in\mathcal F_c$, then
$\sum_{k=1}^K f_k(x_k)\le c(x_1,\dots,x_K)$ for all $(x_1,\dots,x_K)$; integrating
against the product measure $\mu_1\otimes\cdots\otimes\mu_K$ gives
\[
J(\boldsymbol f)=\sum_{k=1}^K\int_{X_k}f_k\,d\mu_k
\le\int_{X_1\times\cdots\times X_K}c\,d(\mu_1\otimes\cdots\otimes\mu_K)\le M,
\]
so $V_c\le M<\infty$. Moreover the constant family $(-M,0,\dots,0)$ is admissible
(since $-M\le c$), whence $V_c\ge -M>-\infty$. Thus $V_c\in\mathbb R$.

\medskip
\noindent\emph{The $c$-conjugate improvement operators.}
For a family $\boldsymbol h=(h_1,\dots,h_K)$ with $h_k:X_k\to\mathbb R$ and an index
$k\in\{1,\dots,K\}$, define $T_k\boldsymbol h$ by
\[
(T_k\boldsymbol h)_k(x_k)
=\inf_{(x_l)_{l\ne k}\in\prod_{l\ne k}X_l}
\Bigl\{c(x_1,\dots,x_K)-\sum_{l\ne k}h_l(x_l)\Bigr\},
\qquad
(T_k\boldsymbol h)_l=h_l\quad(l\ne k).
\]
We record three properties used repeatedly.

\emph{(P1) Monotonicity.} If $\boldsymbol h\le\boldsymbol q$ componentwise, then for
every $(x_l)_{l\ne k}$ we have
$c-\sum_{l\ne k}h_l\ge c-\sum_{l\ne k}q_l$, and taking the infimum gives
$(T_k\boldsymbol h)_k\ge(T_k\boldsymbol q)_k$ — note the order reverses in the
infimum but $T_k$ only changes the $k$-th component while keeping
$(T_k\boldsymbol h)_l=h_l\le q_l=(T_k\boldsymbol q)_l$; hence as full families
$T_k\boldsymbol h$ and $T_k\boldsymbol q$ remain comparable in the sense needed
below, and on the $k$-th component the value increases when the \emph{other}
components decrease. (We use (P1) only through the explicit chain estimate in
Step~1, so no ambiguity arises.)

\emph{(P2) Extensivity and admissibility preservation.} Suppose $\boldsymbol h$ is
admissible, i.e.\ $\sum_{l=1}^K h_l\le c$. Fixing $x_k$ and rearranging,
$h_k(x_k)\le c(x_1,\dots,x_K)-\sum_{l\ne k}h_l(x_l)$ for every $(x_l)_{l\ne k}$;
taking the infimum over $(x_l)_{l\ne k}$ yields
\[
h_k(x_k)\le (T_k\boldsymbol h)_k(x_k),\qquad\forall x_k\in X_k.
\]
Conversely, by definition of the infimum,
$(T_k\boldsymbol h)_k(x_k)\le c(x_1,\dots,x_K)-\sum_{l\ne k}h_l(x_l)$ for every
$(x_1,\dots,x_K)$, i.e.\
$(T_k\boldsymbol h)_k(x_k)+\sum_{l\ne k}h_l(x_l)\le c(x_1,\dots,x_K)$. Since the
other components are unchanged, $T_k\boldsymbol h$ is admissible.

\emph{(P3) Fixed points.} By Definition~\ref{con}, $\boldsymbol h$ is mutually
$c$-conjugate if and only if $h_k=(T_k\boldsymbol h)_k$ on $X_k$ for every $k$,
i.e.\ if and only if $T_k\boldsymbol h=\boldsymbol h$ for all $k$.

\medskip
\noindent\emph{Step 1: construction of a maximizing sequence in $\mathcal G_c$.}
By the compact Kantorovich duality theorem~\ref{KANTO},
\[
V_c=\sup\Bigl\{J(\boldsymbol u):\ \boldsymbol u\in\mathcal F_c\cap\textstyle\prod_{k=1}^K C(X_k)\Bigr\},
\]
so there is a maximizing sequence
$\boldsymbol u^n=(u_1^n,\dots,u_K^n)\in\mathcal F_c\cap\prod_k C(X_k)$ with
$J(\boldsymbol u^n)\to V_c$. Fix $n\ge1$ and introduce the class
\[
\mathcal A_n=\Bigl\{\boldsymbol h=(h_1,\dots,h_K):\ h_k:X_k\to\mathbb R,\
h_k\ge u_k^n\ \text{on }X_k,\ \sum_{k=1}^K h_k\le c\Bigr\},
\]
partially ordered by the componentwise order
$\boldsymbol h\le\boldsymbol q\iff h_k\le q_k$ on $X_k$ for all $k$. The set
$\mathcal A_n$ is nonempty since $\boldsymbol u^n\in\mathcal A_n$.

\emph{Every chain has an upper bound.} Let $\mathcal C\subset\mathcal A_n$ be a
totally ordered chain. For each $k$ define
\[
h_k^*(x_k)=\sup_{\boldsymbol h\in\mathcal C}h_k(x_k).
\]
We first check $h_k^*$ is real-valued. Fix reference points $\bar x_l\in X_l$ for
$l\ne k$. For every $\boldsymbol h\in\mathcal C$, admissibility at
$(\,\cdot\,,\bar x_l)_{l\ne k}$ gives
$h_k(x_k)+\sum_{l\ne k}h_l(\bar x_l)\le c$, and since $h_l\ge u_l^n$,
\[
h_k(x_k)\le M-\sum_{l\ne k}u_l^n(\bar x_l)<\infty,
\]
a bound independent of $\boldsymbol h\in\mathcal C$. Hence
$h_k^*(x_k)<\infty$; and $h_k^*\ge u_k^n>-\infty$. Thus $h_k^*$ is real-valued and
$h_k^*\ge u_k^n$.

We verify $\boldsymbol h^*=(h_1^*,\dots,h_K^*)$ is admissible. Fix
$(x_1,\dots,x_K)\in X_1\times\cdots\times X_K$ and $\eta>0$. For each
$k=1,\dots,K$ choose $\boldsymbol h^{(k)}\in\mathcal C$ with
\[
h_k^*(x_k)\le h_k^{(k)}(x_k)+\frac{\eta}{K}.
\]
Because $\mathcal C$ is totally ordered and $\{\boldsymbol h^{(1)},\dots,\boldsymbol h^{(K)}\}$
is finite, one of these elements, say $\boldsymbol h^{(0)}\in\mathcal C$, dominates
all the others, so $h_k^{(k)}(x_k)\le h_k^{(0)}(x_k)$ for every $k$. Therefore
\[
\sum_{k=1}^K h_k^*(x_k)
\le\sum_{k=1}^K\Bigl(h_k^{(0)}(x_k)+\frac{\eta}{K}\Bigr)
=\sum_{k=1}^K h_k^{(0)}(x_k)+\eta
\le c(x_1,\dots,x_K)+\eta,
\]
using admissibility of $\boldsymbol h^{(0)}\in\mathcal A_n$. Letting
$\eta\downarrow0$ gives $\sum_k h_k^*(x_k)\le c(x_1,\dots,x_K)$. Hence
$\boldsymbol h^*\in\mathcal A_n$, and $h_k\le h_k^*$ for all
$\boldsymbol h\in\mathcal C$ and all $k$ shows $\boldsymbol h^*$ is an upper bound
of $\mathcal C$.
By Zorn's lemma (Lemma~\ref{Zorn}), $\mathcal A_n$ contains a maximal element
$\boldsymbol g^n=(g_1^n,\dots,g_K^n)$.

\emph{$\boldsymbol g^n$ is mutually $c$-conjugate.} Fix $k$. Since
$\boldsymbol g^n$ is admissible, (P2) gives $T_k\boldsymbol g^n\in\mathcal A_n$
(its components still dominate $u^n$, because $(T_k\boldsymbol g^n)_k\ge g_k^n\ge u_k^n$
and the other components are unchanged) and $T_k\boldsymbol g^n\ge\boldsymbol g^n$.
Also $T_k\boldsymbol g^n$ is real-valued: it is bounded below by $g_k^n$, and
bounded above by
\[
(T_k\boldsymbol g^n)_k(x_k)\le c(\bar x_1,\dots,\bar x_{k-1},x_k,\bar x_{k+1},\dots,\bar x_K)
-\sum_{l\ne k}g_l^n(\bar x_l)<\infty
\]
for fixed reference points $\bar x_l$. By maximality of $\boldsymbol g^n$ we must
have $T_k\boldsymbol g^n=\boldsymbol g^n$, i.e.\ $g_k^n=(T_k\boldsymbol g^n)_k$ on
$X_k$. As $k$ was arbitrary, (P3) shows $\boldsymbol g^n$ is mutually
$c$-conjugate.

\emph{$\boldsymbol g^n\in\mathcal G_c$.} Each $g_k^n$ is a real-valued partial
$c$-transform, so by Lemma~\ref{equi} it inherits a modulus of continuity of $c$;
in particular $g_k^n$ is continuous, hence measurable, and bounded (as $X_k$ is
compact), so $\int_{X_k}|g_k^n|\,d\mu_k<\infty$ and $g_k^n\in L_1(\mu_k)$.
Together with admissibility, $\boldsymbol g^n\in\mathcal F_c$, and with mutual
$c$-conjugacy, $\boldsymbol g^n\in\mathcal G_c$.

\emph{$\boldsymbol g^n$ is maximizing.} Since $g_k^n\ge u_k^n$ for every $k$ and
each $\mu_k$ is a probability measure,
$J(\boldsymbol g^n)\ge J(\boldsymbol u^n)$; and $\boldsymbol g^n\in\mathcal F_c$
gives $J(\boldsymbol g^n)\le V_c$. From
$J(\boldsymbol u^n)\le J(\boldsymbol g^n)\le V_c$ and
$J(\boldsymbol u^n)\to V_c$, the squeeze theorem yields
$J(\boldsymbol g^n)\to V_c$. Thus
$(\boldsymbol g^n)_{n\ge1}\subset\mathcal G_c$ is a maximizing sequence for the
dual problem over $\mathcal F_c$.

\medskip
\noindent\emph{Step 2: equicontinuity, normalization, and uniform bounds.}
By Lemma~\ref{equi}, for each $k$ the family $(g_k^n)_{n\ge1}$ is equicontinuous
on $X_k$ with a modulus depending only on $c$ (independent of $n$). By
Lemma~\ref{uniform-bounds}, there exist constants $a_1^n,\dots,a_K^n$ with
$\sum_{k=1}^K a_k^n=0$ such that the translated families
$\widetilde{\boldsymbol g}^n=(g_1^n+a_1^n,\dots,g_K^n+a_K^n)$ still belong to
$\mathcal G_c$, still satisfy $J(\widetilde{\boldsymbol g}^n)=J(\boldsymbol g^n)$
(because the added constants sum to zero and each $\mu_k$ is a probability
measure), and obey uniform bounds $\|g_k^n+a_k^n\|_\infty\le C_k<\infty$, with
$C_k$ independent of $n$. Relabeling $\widetilde{\boldsymbol g}^n$ as
$\boldsymbol g^n$ (which changes neither $\mathcal G_c$-membership nor the value of
$J$ nor the equicontinuity), we may henceforth assume
\[
\|g_k^n\|_\infty\le C_k\qquad\text{for all }n,\ k=1,\dots,K.
\]

\medskip
\noindent\emph{Step 3: passage to the limit via Arzelà--Ascoli.}
For each $k$ the sequence $(g_k^n)_{n}$ is uniformly bounded and equicontinuous on
the compact metric space $X_k$. By the Arzelà--Ascoli theorem
(Theorem~\ref{arz}), passing to a subsequence (the same subsequence for all $k$,
by extracting successively over $k=1,\dots,K$), there exist continuous functions
$f_k:X_k\to\mathbb R$ with
\[
g_k^n\longrightarrow f_k\quad\text{uniformly on }X_k,\qquad k=1,\dots,K.
\]
Write $\boldsymbol f=(f_1,\dots,f_K)$.

\emph{Admissibility of $\boldsymbol f$.} From $\sum_k g_k^n\le c$, uniform
convergence gives $\sum_k f_k\le c$, so $\boldsymbol f\in\mathcal F_c$.

\emph{Optimality of $\boldsymbol f$.} Uniform convergence implies
$\int_{X_k}g_k^n\,d\mu_k\to\int_{X_k}f_k\,d\mu_k$ for each $k$, hence
$J(\boldsymbol g^n)\to J(\boldsymbol f)$. Since $J(\boldsymbol g^n)\to V_c$, we get
$J(\boldsymbol f)=V_c=\sup_{\boldsymbol h\in\mathcal F_c}J(\boldsymbol h)$, so
$\boldsymbol f$ maximizes the dual problem over $\mathcal F_c$.

\emph{$\boldsymbol f\in\mathcal G_c$.} Fix $k$. Since
$\boldsymbol g^n\in\mathcal G_c$, for every $x_k\in X_k$,
\[
g_k^n(x_k)=\inf_{(x_l)_{l\ne k}}\Bigl\{c(x_1,\dots,x_K)-\sum_{l\ne k}g_l^n(x_l)\Bigr\}
=\inf_{(x_l)_{l\ne k}}H_n(x_1,\dots,x_K),
\]
with $H_n=c-\sum_{l\ne k}g_l^n$ and $H=c-\sum_{l\ne k}f_l$. Since $g_l^n\to f_l$
uniformly on $X_l$ for every $l\ne k$, we have $H_n\to H$ uniformly on the compact
product. The infimum operation is $1$-Lipschitz for the supremum norm: for any
bounded $\varphi,\psi$,
$|\inf\varphi-\inf\psi|\le\|\varphi-\psi\|_\infty$; applying this with the infimum
taken over $(x_l)_{l\ne k}$ at each fixed $x_k$,
\[
\Bigl\|\inf_{(x_l)_{l\ne k}}H_n-\inf_{(x_l)_{l\ne k}}H\Bigr\|_{\infty,\,x_k}
\le\|H_n-H\|_\infty\longrightarrow0.
\]
Thus $g_k^n=\inf_{(x_l)_{l\ne k}}H_n\to\inf_{(x_l)_{l\ne k}}H$ uniformly in $x_k$;
since also $g_k^n\to f_k$ uniformly, the limits coincide:
\[
f_k(x_k)=\inf_{(x_l)_{l\ne k}}\Bigl\{c(x_1,\dots,x_K)-\sum_{l\ne k}f_l(x_l)\Bigr\},
\qquad\forall x_k\in X_k.
\]
This holds for every $k$, so $T_k\boldsymbol f=\boldsymbol f$ for all $k$, and by
(P3), $\boldsymbol f\in\mathcal G_c$.

\medskip
\noindent\emph{Conclusion.}
The maximizer $\boldsymbol f$ lies in $\mathcal G_c\subset\mathcal F_c$ and attains
the supremum over $\mathcal F_c$. Hence, on the one hand
$\max_{\boldsymbol h\in\mathcal G_c}J(\boldsymbol h)\le
\max_{\boldsymbol h\in\mathcal F_c}J(\boldsymbol h)$ (as $\mathcal G_c\subset\mathcal F_c$),
and on the other hand $\boldsymbol f\in\mathcal G_c$ attains the right-hand side, so
\[
\max_{\boldsymbol h\in\mathcal G_c}J(\boldsymbol h)
=\max_{\boldsymbol h\in\mathcal F_c}J(\boldsymbol h)=V_c.
\]
Finally, by Theorem~\ref{KANTO}, $\mathsf{MOT}_c(\mu_1,\dots,\mu_K)=V_c$, whence
\[
\mathsf{MOT}_c(\mu_1,\dots,\mu_K)
=\max_{\boldsymbol f\in\mathcal F_c}J(\boldsymbol f)
=\max_{\boldsymbol f\in\mathcal G_c}J(\boldsymbol f).
\]
This completes the proof.
\hfill\(\Box\)



\subsubsection*{Proof of Theorem~\ref{noncompact}}

Throughout, set $M=\|c\|_\infty<\infty$, which is finite because $c$ is bounded.
Write $E=X_1\times\cdots\times X_K$ and abbreviate a generic point by
$\boldsymbol x=(x_1,\dots,x_K)$.

\medskip
\noindent\emph{Step 0: the hypothesis and what must be shown.}
By assumption, there exists an optimal plan
$\pi^*\in\Pi(\mu_1,\dots,\mu_K)$, optimality meaning
$I(\pi^*)=\min_{\pi\in\Pi(\mu_1,\dots,\mu_K)}\int_E c\,d\pi$, the minimum being
attained by Lemma~\ref{compact} together with the narrow continuity of
$\pi\mapsto\int c\,d\pi$ for $c\in\mathcal C_b(E)$, such that
\[
\Gamma=\operatorname{supp}(\pi^*)
\]
admits a real-valued, Borel measurable, mutually $c$-conjugate splitting family
$\boldsymbol g=(g_1,\dots,g_K)$, $g_k:X_k\to\mathbb R$. By Definition of a
splitting family, $\boldsymbol g$ satisfies
\begin{align}
\sum_{k=1}^K g_k(x_k)&\le c(\boldsymbol x), &&\forall\,\boldsymbol x\in E,
\label{eq:nc-admiss}\\
\sum_{k=1}^K g_k(x_k)&= c(\boldsymbol x), &&\forall\,\boldsymbol x\in\Gamma,
\label{eq:nc-satur}
\end{align}
and, being mutually $c$-conjugate (Definition~\ref{con}), for every
$k\in\{1,\dots,K\}$ and every $x_k\in X_k$,
\begin{equation}\label{eq:nc-conj}
g_k(x_k)=\inf_{(x_l)_{l\ne k}\in\prod_{l\ne k}X_l}
\Bigl\{c(\boldsymbol x)-\sum_{l\ne k}g_l(x_l)\Bigr\}.
\end{equation}
We must show that $\boldsymbol g$ is an admissible dual family
($\boldsymbol g\in\mathcal F_c$, indeed $\boldsymbol g\in\mathcal G_c$), that it
attains the dual supremum, and that the three values coincide with the primal
minimum.

\medskip
\noindent\emph{Step 1: boundedness and integrability of the potentials.}
By Proposition~\ref{bounded-splitting}, since $c$ is bounded and
$\boldsymbol g$ is a real-valued mutually $c$-conjugate splitting family, each
$g_k$ is bounded; explicitly, fixing any
$\bar{\boldsymbol x}=(\bar x_1,\dots,\bar x_K)\in\Gamma$ (which is nonempty as
$\pi^*$ is a probability measure, so $\Gamma\ne\varnothing$), the conjugacy
identity~\eqref{eq:nc-conj} evaluated at $(x_l)_{l\ne k}=(\bar x_l)_{l\ne k}$ gives
\[
g_k(x_k)\le c(\bar x_1,\dots,\bar x_{k-1},x_k,\bar x_{k+1},\dots,\bar x_K)
-\sum_{l\ne k}g_l(\bar x_l)\le M-\sum_{l\ne k}g_l(\bar x_l)=:B_k<\infty,
\]
a finite upper bound; the corresponding lower bound
$g_k\ge -M-\sum_{l\ne k}B_l$ then follows as in
Proposition~\ref{bounded-splitting}. Hence there is a finite constant $B$ with
$\|g_k\|_\infty\le B$ for all $k$. Since each $g_k$ is Borel measurable and
bounded, and each $\mu_k$ is a probability measure,
\[
\int_{X_k}|g_k|\,d\mu_k\le B<\infty,\qquad\text{so}\qquad g_k\in L_1(\mu_k),
\quad k=1,\dots,K.
\]
Combined with the admissibility~\eqref{eq:nc-admiss}, this yields
$\boldsymbol g\in\mathcal F_c$; combined further with the conjugacy
identities~\eqref{eq:nc-conj}, $\boldsymbol g\in\mathcal G_c$.

\medskip
\noindent\emph{Step 2: the dual value at $\boldsymbol g$ equals the primal
minimum.}
Because $\pi^*\in\Pi(\mu_1,\dots,\mu_K)$, its $k$-th marginal is $\mu_k$, so for
each $k$, $\int_{X_k}g_k\,d\mu_k=\int_E g_k(x_k)\,d\pi^*(\boldsymbol x)$; summing,
\begin{equation}\label{eq:nc-Jg}
J(\boldsymbol g)
=\sum_{k=1}^K\int_{X_k}g_k\,d\mu_k
=\int_E\Bigl(\sum_{k=1}^K g_k(x_k)\Bigr)\,d\pi^*(\boldsymbol x).
\end{equation}
We now show the integrand coincides with $c$ for $\pi^*$-a.e.\ $\boldsymbol x$.
The function $\Phi(\boldsymbol x)=c(\boldsymbol x)-\sum_{k=1}^K g_k(x_k)$ is
Borel measurable, nonnegative by~\eqref{eq:nc-admiss}, and vanishes on
$\Gamma=\operatorname{supp}(\pi^*)$ by~\eqref{eq:nc-satur}. Since $\pi^*$ is, by
definition of the support, concentrated on its support — that is,
$\pi^*\bigl(E\setminus\Gamma\bigr)=0$ — we have $\Phi=0$ $\pi^*$-almost
everywhere, whence
\[
\int_E\Phi\,d\pi^*=0,
\qquad\text{i.e.}\qquad
\int_E\Bigl(\sum_{k=1}^K g_k(x_k)\Bigr)\,d\pi^*
=\int_E c\,d\pi^*=I(\pi^*).
\]
(The identity $\pi^*(E\setminus\Gamma)=0$ holds because $E$ is Polish, hence
second countable, so the support of a Borel probability measure carries full
mass.) Substituting into~\eqref{eq:nc-Jg} and using optimality of $\pi^*$,
\begin{equation}\label{eq:nc-Jg-min}
J(\boldsymbol g)=I(\pi^*)
=\min_{\pi\in\Pi(\mu_1,\dots,\mu_K)}\int_E c\,d\pi.
\end{equation}

\medskip
\noindent\emph{Step 3: weak duality.}
For every $\boldsymbol f=(f_1,\dots,f_K)\in\mathcal F_c$ and every
$\pi\in\Pi(\mu_1,\dots,\mu_K)$, admissibility
$\sum_k f_k(x_k)\le c(\boldsymbol x)$ and the marginal constraints give
\[
J(\boldsymbol f)=\sum_{k=1}^K\int_{X_k}f_k\,d\mu_k
=\int_E\Bigl(\sum_{k=1}^K f_k(x_k)\Bigr)\,d\pi
\le\int_E c\,d\pi=I(\pi).
\]
Taking the supremum over $\boldsymbol f\in\mathcal F_c$ and the infimum over
$\pi$,
\begin{equation}\label{eq:nc-weak}
\sup_{\boldsymbol f\in\mathcal F_c}J(\boldsymbol f)
\le\min_{\pi\in\Pi(\mu_1,\dots,\mu_K)}\int_E c\,d\pi.
\end{equation}
(This is the content of Lemma~\ref{easyy}, restated here for the bounded
continuous cost and the attained minimum.)

\medskip
\noindent\emph{Step 4: dual attainment over $\mathcal F_c$.}
Since $\boldsymbol g\in\mathcal F_c$, the supremum dominates $J(\boldsymbol g)$:
\[
\sup_{\boldsymbol f\in\mathcal F_c}J(\boldsymbol f)\ge J(\boldsymbol g).
\]
On the other hand, by~\eqref{eq:nc-Jg-min} and the weak
duality~\eqref{eq:nc-weak},
\[
J(\boldsymbol g)
=\min_{\pi\in\Pi(\mu_1,\dots,\mu_K)}\int_E c\,d\pi
\ge\sup_{\boldsymbol f\in\mathcal F_c}J(\boldsymbol f).
\]
The two displays force equality:
\begin{equation}\label{eq:nc-attain-Fc}
J(\boldsymbol g)=\sup_{\boldsymbol f\in\mathcal F_c}J(\boldsymbol f)
=\max_{\boldsymbol f\in\mathcal F_c}J(\boldsymbol f),
\end{equation}
the supremum being attained precisely at $\boldsymbol g$.

\medskip
\noindent\emph{Step 5: attainment over the canonical class $\mathcal G_c$.}
Since $\mathcal G_c\subset\mathcal F_c$, restricting the supremum can only
decrease it:
\[
\sup_{\boldsymbol f\in\mathcal G_c}J(\boldsymbol f)
\le\sup_{\boldsymbol f\in\mathcal F_c}J(\boldsymbol f).
\]
Conversely, $\boldsymbol g\in\mathcal G_c$ (Step~1), so by~\eqref{eq:nc-attain-Fc},
\[
\sup_{\boldsymbol f\in\mathcal G_c}J(\boldsymbol f)
\ge J(\boldsymbol g)
=\sup_{\boldsymbol f\in\mathcal F_c}J(\boldsymbol f).
\]
The two inequalities give
\begin{equation}\label{eq:nc-attain-Gc}
\sup_{\boldsymbol f\in\mathcal G_c}J(\boldsymbol f)
=\sup_{\boldsymbol f\in\mathcal F_c}J(\boldsymbol f)
=J(\boldsymbol g),
\end{equation}
and the supremum over $\mathcal G_c$ is attained at $\boldsymbol g$; in particular
$\max_{\boldsymbol f\in\mathcal F_c}J(\boldsymbol f)=\max_{\boldsymbol f\in\mathcal G_c}J(\boldsymbol f)$.

\medskip
\noindent\emph{Conclusion.}
Combining~\eqref{eq:nc-Jg-min}, \eqref{eq:nc-attain-Fc}
and~\eqref{eq:nc-attain-Gc}, and recalling
$\mathsf{MOT}_c(\mu_1,\dots,\mu_K)=\min_{\pi\in\Pi(\mu_1,\dots,\mu_K)}\int_E c\,d\pi$,
\[
\mathsf{MOT}_c(\mu_1,\dots,\mu_K)
=\min_{\pi\in\Pi(\mu_1,\dots,\mu_K)}\int_{E} c\,d\pi
=\max_{\boldsymbol f\in\mathcal F_c}J(\boldsymbol f)
=\max_{\boldsymbol f\in\mathcal G_c}J(\boldsymbol f),
\]
with all three optima attained, the dual ones at the bounded Borel measurable
mutually $c$-conjugate family $\boldsymbol g\in\mathcal G_c$. This completes the
proof.
\hfill\(\Box\)















\numberwithin{theorem}{section}
\numberwithin{definition}{section}
\numberwithin{lemma}{section}
\subsection{Supporting results}


\begin{theorem}[Fenchel--Rockafellar duality \cite{villani2003}]\label{fenchell}
Let $E$ be a normed vector space, let $E^*$ be its topological dual space, and let 
$\Theta, \Xi : E \to \mathbb{R} \cup \{+\infty\}$ be two convex functions.
Denote by $\Theta^*$ and $\Xi^*$ their respective Legendre-Fenchel transforms.
Suppose that there exists $z_0 \in E$ such that 
\begin{align*}
    \Theta(z_0) < \infty, \quad \Xi(z_0) < \infty,
    \quad 
\end{align*}
and $\Theta \text{ is continuous at } z_0$. Then,
\begin{equation*}
    \inf_{z \in E} \big[ \Theta(z) + \Xi(z) \big]
    \;=\;
    \max_{z^* \in E^*} \big[ -\Theta^*(-z^*) - \Xi^*(z^*) \big].
\end{equation*}
\end{theorem}


\begin{theorem}[Riesz representation theorem \cite{knapp2007}]\label{RIEZ}
Let $X_1, \dots, X_K$ be compact metric spaces and denote by $ E=\mathcal{C}(X_1 \times\ldots \times X_K)$ the space of continuous and real-valued functions
on the product, equipped with the supremum norm. 
Then, the dual space $E^*$ can be identified with the space $\mathcal{M}\!(X_1 \times\ldots \times X_K)$ of finite signed regular Borel measures, equipped with the total variation norm. The correspondence is given by
\begin{align*}
\mu \in \mathcal M(X_1 \times \cdots \times X_K)
\;\longmapsto\;
\mathcal{L}_\mu \in E^{*},
\end{align*}
where, for every $u \in E$,
\begin{align*}
\mathcal{L}_\mu(u) 
= \int_{X_1 \times\ldots \times X_K} 
u(x_1, \dots, x_K)\, d\mu(x_1, \dots, x_K).
\end{align*}
\end{theorem}


\begin{theorem}
[Arzelà-Ascoli \cite{santambrogio2015}]\label{arz}
Let $X$ be a compact metric space, and let $(f_n)_{n\in\mathbb{N}}$ 
be a sequence of functions $f_n : X \to \mathbb{R}$.
Assume that $(f_n)$ is equicontinuous \footnote{A family of functions $(f_n)_{n \in \mathbb{N}} \subset \mathcal{C}(X,\mathbb{R})$ is \textit{equicontinuous} if, 
for every $\varepsilon > 0$, there exists $\delta > 0$ such that for all $x, y \in X$ and all $n \in \mathbb{N}$, $|x - y| < \delta \ \Rightarrow \ |f_n(x) - f_n(y)| < \varepsilon$.}, and uniformly bounded. Then, there exists a subsequence $(f_{n_k})_{k\in\mathbb{N}}$ 
that converges uniformly on $X$ to a continuous function 
$f : X \to \mathbb{R}$.
\end{theorem}



\begin{theorem}[Prokhorov's Theorem \cite{villani2008}]\label{prokhorov}
Let $X$ be a Polish space. A subset $P \subset \mathcal{P}(X)$ is precompact\footnote{A subset of a topological space is called precompact (or relatively compact) if its closure is compact.}
for the weak topology if and only if it is tight.
\end{theorem}


\begin{definition}\label{tightness}
(Tightness \cite{villani2008})
Let $X$ be a Polish space. A family $P \subset \mathcal{P}(X)$ is said to be \emph{tight} if, 
for every $\varepsilon>0$, there exists a compact set $\mathcal{K}_\varepsilon \subset X$ such that
\begin{align*}
\mu[X \setminus \mathcal{K}_\varepsilon] \le \varepsilon,
\qquad \forall\, \mu \in P.
\end{align*}
\end{definition}

\begin{definition}[Weak convergence \cite{van2013}]\label{cv}
Let $X_n : \Omega \to \mathcal{X}$ and $X : \Omega \to \mathcal{X}$ be arbitrary measurable maps. If $X_n$ and $X$ are $\mathcal{X}$-valued random variables with respective laws 
$\mathbb{P}_n$ and $\mathbb{P}$, then
\begin{align*}
X_n \rightsquigarrow X
\quad \Longleftrightarrow \quad
\int_{\mathcal{X}} f\, d\mathbb{P}_n 
\longrightarrow 
\int_{\mathcal{X}} f\, d\mathbb{P},
\qquad \text{for all } f \in \mathcal{C}_b(\mathcal{X}).
\end{align*}
\end{definition}


\begin{lemma}[\cite{knapp2007}]
Let $(X, \mathcal{B}(X))$ be a Polish space, and let $\mu$ be a Borel probability measure on $X$. Then $\mu$ is regular: for every Borel set $A \subset X$,
\begin{align*}
\mu(A) 
&= \sup \{\, \mu(K) : K \subset A,\, K \text{ compact} \,\}
\\
&= \inf \{\, \mu(O) : A \subset O,\, O \text{ open} \,\}.
\end{align*}
\end{lemma}


\begin{lemma}[Zorn's lemma \cite{Paul1960}]\label{Zorn}
Let \(X\) be a partially ordered set. If every chain in \(X\) has an upper
bound in \(X\), then \(X\) contains a maximal element.
\end{lemma}


\begin{thebibliography}{10}

\bibitem{Beiglb}
Mathias Beiglb\"{o}ck, Christian L\'{e}onard, and Walter Schachermayer.
\newblock A generalized dual maximizer for the {M}onge-{K}antorovich transport
  problem.
\newblock {\em ESAIM Probab. Stat.}, 16:306--323, 2012.

\bibitem{Bogachev}
V.~I. Bogachev and A.~V. Kolesnikov.
\newblock The {M}onge-{K}antorovich problem: achievements, connections, and
  prospects.
\newblock {\em Uspekhi Mat. Nauk}, 67(5(407)):3--110, 2012.

\bibitem{bourbaki1998}
Nicolas Bourbaki.
\newblock {\em General topology. {C}hapters 5--10}.
\newblock Elements of Mathematics (Berlin). Springer-Verlag, Berlin, 1998.
\newblock Translated from the French, Reprint of the 1989 English translation.

\bibitem{cao2019}
Jiezhang Cao, Langyuan Mo, Yifan Zhang, Kui Jia, Chunhua Shen, and Mingkui Tan.
\newblock Multi-marginal wasserstein gan.
\newblock {\em Advances in Neural Information Processing Systems}, 32, 2019.

\bibitem{carlier2010}
G.~Carlier and I.~Ekeland.
\newblock Matching for teams.
\newblock {\em Econom. Theory}, 42(2):397--418, 2010.

\bibitem{colombo2015}
Maria Colombo, Luigi De~Pascale, and Simone Di~Marino.
\newblock Multimarginal optimal transport maps for one-dimensional repulsive
  costs.
\newblock {\em Canad. J. Math.}, 67(2):350--368, 2015.

\bibitem{del2019}
Eustasio del Barrio and Jean-Michel Loubes.
\newblock Central limit theorems for empirical transportation cost in general
  dimension.
\newblock {\em Ann. Probab.}, 47(2):926--951, 2019.

\bibitem{deshpande2018}
Ishan Deshpande, Ziyu Zhang, and Alexander~G Schwing.
\newblock Generative modeling using the sliced wasserstein distance.
\newblock In {\em Proceedings of the IEEE conference on computer vision and
  pattern recognition}, pages 3483--3491, 2018.

\bibitem{fournier2015}
Nicolas Fournier and Arnaud Guillin.
\newblock On the rate of convergence in {W}asserstein distance of the empirical
  measure.
\newblock {\em Probab. Theory Related Fields}, 162(3-4):707--738, 2015.

\bibitem{gangbo1998}
Wilfrid Gangbo and Andrzej \'{S}wi\c{e}ch.
\newblock Optimal maps for the multidimensional {M}onge-{K}antorovich problem.
\newblock {\em Comm. Pure Appl. Math.}, 51(1):23--45, 1998.

\bibitem{genevay2018}
Aude Genevay, Gabriel Peyr{\'e}, and Marco Cuturi.
\newblock Learning generative models with sinkhorn divergences.
\newblock In {\em International Conference on Artificial Intelligence and
  Statistics}, pages 1608--1617. PMLR, 2018.

\bibitem{griessler2016}
Claus Griessler.
\newblock {$C$}-cyclical monotonicity as a sufficient criterion for optimality
  in the multimarginal {M}onge-{K}antorovich problem.
\newblock {\em Proc. Amer. Math. Soc.}, 146(11):4735--4740, 2018.

\bibitem{Paul1960}
Paul~R. Halmos.
\newblock {\em Naive set theory}.
\newblock Undergraduate Texts in Mathematics. Springer-Verlag, New
  York-Heidelberg, 1960 edition, 1974.

\bibitem{He2019}
Zhenliang He, Wangmeng Zuo, Meina Kan, Shiguang Shan, and Xilin Chen.
\newblock Attgan: Facial attribute editing by only changing what you want.
\newblock {\em IEEE transactions on image processing}, 28(11):5464--5478, 2019.

\bibitem{heinich2002}
Henri Heinich.
\newblock Probl\`eme de {M}onge pour {$n$} probabilit\'{e}s.
\newblock {\em C. R. Math. Acad. Sci. Paris}, 334(9):793--795, 2002.

\bibitem{Hui2018}
Le~Hui, Xiang Li, Jiaxin Chen, Hongliang He, and Jian Yang.
\newblock Unsupervised multi-domain image translation with domain-specific
  encoders/decoders.
\newblock In {\em 2018 24th International Conference on Pattern Recognition
  (ICPR)}, pages 2044--2049. IEEE, 2018.

\bibitem{article}
L.~Kantorovitch.
\newblock On the translocation of masses.
\newblock {\em C. R. (Doklady) Acad. Sci. URSS (N.S.)}, 37:199--201, 1942.

\bibitem{kellerer}
Hans~G. Kellerer.
\newblock Duality theorems for marginal problems.
\newblock {\em Z. Wahrsch. Verw. Gebiete}, 67(4):399--432, 1984.

\bibitem{knapp2007}
Anthony~W. Knapp.
\newblock {\em Basic real analysis}.
\newblock Cornerstones. Birkh\"{a}user Boston, Inc., Boston, MA, 2005.
\newblock Along with a companion volume {\it Advanced real analysis}.

\bibitem{knott1984}
M.~Knott and C.~S. Smith.
\newblock On the optimal mapping of distributions.
\newblock {\em J. Optim. Theory Appl.}, 43(1):39--49, 1984.

\bibitem{Krylov2002}
N.~V. Krylov.
\newblock {\em Introduction to the theory of random processes}, volume~43 of
  {\em Graduate Studies in Mathematics}.
\newblock American Mathematical Society, Providence, RI, 2002.

\bibitem{monge1781memoire}
G.~Monge.
\newblock {\em M{\'e}moire sur la th{\'e}orie des d{\'e}blais et des remblais}.
\newblock Imprimerie royale, 1781.

\bibitem{pass2012}
Brendan Pass.
\newblock On the local structure of optimal measures in the multi-marginal
  optimal transportation problem.
\newblock {\em Calc. Var. Partial Differential Equations}, 43(3-4):529--536,
  2012.

\bibitem{pass2015}
Brendan Pass.
\newblock Multi-marginal optimal transport: theory and applications.
\newblock {\em ESAIM Math. Model. Numer. Anal.}, 49(6):1771--1790, 2015.

\bibitem{pass2022}
Brendan Pass and Adolfo Vargas-Jim\'{e}nez.
\newblock Monge solutions and uniqueness in multi-marginal optimal transport
  via graph theory.
\newblock {\em Adv. Math.}, 428:Paper No. 109101, 43, 2023.

\bibitem{peyre2019}
Gabriel Peyr{\'e}, Marco Cuturi, et~al.
\newblock Computational optimal transport: With applications to data science.
\newblock {\em Foundations and Trends{\textregistered} in Machine Learning},
  11(5-6):355--607, 2019.

\bibitem{rachev1998}
Svetlozar~T. Rachev and Ludger R\"{u}schendorf.
\newblock {\em Mass transportation problems. {V}ol. {I}}.
\newblock Probability and its Applications (New York). Springer-Verlag, New
  York, 1998.
\newblock Theory.

\bibitem{Ramachandran}
D.~Ramachandran and L.~R\"{u}schendorf.
\newblock A general duality theorem for marginal problems.
\newblock {\em Probab. Theory Related Fields}, 101(3):311--319, 1995.

\bibitem{Rigo}
Pietro Rigo.
\newblock A note on duality theorems in mass transportation.
\newblock {\em J. Theoret. Probab.}, 33(4):2337--2350, 2020.

\bibitem{Rigo2}
Pietro Rigo.
\newblock Finitely additive mass transportation.
\newblock {\em Bernoulli}, 30(3):1825--1844, 2024.

\bibitem{santambrogio2015}
Filippo Santambrogio.
\newblock {\em Optimal transport for applied mathematicians}, volume~87 of {\em
  Progress in Nonlinear Differential Equations and their Applications}.
\newblock Birkh\"{a}user/Springer, Cham, 2015.
\newblock Calculus of variations, PDEs, and modeling.

\bibitem{Sudakov1979}
V.~N. Sudakov.
\newblock Geometric problems in the theory of infinite-dimensional probability
  distributions.
\newblock {\em Proc. Steklov Inst. Math.}, (2):i--v, 1--178, 1979.

\bibitem{tameling2019}
Carla Tameling, Max Sommerfeld, and Axel Munk.
\newblock Empirical optimal transport on countable metric spaces:
  distributional limits and statistical applications.
\newblock {\em Ann. Appl. Probab.}, 29(5):2744--2781, 2019.

\bibitem{van2013}
Aad~W. van~der Vaart and Jon~A. Wellner.
\newblock {\em Weak convergence and empirical processes}.
\newblock Springer Series in Statistics. Springer-Verlag, New York, 1996.
\newblock With applications to statistics.

\bibitem{villani2003}
C\'{e}dric Villani.
\newblock {\em Topics in optimal transportation}, volume~58 of {\em Graduate
  Studies in Mathematics}.
\newblock American Mathematical Society, Providence, RI, 2003.

\bibitem{villani2008}
C\'{e}dric Villani.
\newblock {\em Optimal transport}, volume 338 of {\em Grundlehren der
  mathematischen Wissenschaften [Fundamental Principles of Mathematical
  Sciences]}.
\newblock Springer-Verlag, Berlin, 2009.
\newblock Old and new.

\end{thebibliography}

\bibliographystyle{plain} 


\end{document}
